\chapter{简正模问题}


平衡附近的运动常可用线性方程近似。在单一振动模态中，各坐标以同一频率运动，振幅之比保持不变，这组比例称为振型。稳定平衡附近的任意线性小振动由这些模态叠加而成，初始条件决定各模态的份额。



\section{平衡构形、二次能量与固有频率}

小振动围绕满足约束的平衡构形展开。独立位移确定后，动能与势能的二次项在同一组坐标中给出质量矩阵和刚度矩阵，由此求固有频率。



\subsection{从平衡位置到二次能量}
设独立坐标为 $q_i$，平衡位置为 $q_{i*}$，小偏移为 $\eta_i=q_i-q_{i*}$。以下假定约束不显含时间、真实力可由势能表示，且平衡态静止。若平衡态本身持续转动，还须保留转动参考系中的陀螺项，方程形式会有所不同。

在平衡点附近，将能量保留到小量的二阶：
\begin{equation}
T_2=\frac12\dot{\vct\eta}^{\mathsf T}M\dot{\vct\eta},\qquad
U_2=U_*+\frac12\vct\eta^{\mathsf T}K\vct\eta.
\label{sup4:eq:d7-energy}
\end{equation}
$M$ 是平衡位置处的质量矩阵，$K_{ij}=\partial_i\partial_jU|_*$ 是刚度矩阵，二者均对称。使用独立且非奇异的坐标时，$M$ 正定；$K$ 则由平衡性质决定。

对单坐标系统，若 $T_2=A_*\dot\eta^2/2$、$U_2=U_*+B\eta^2/2$，有
\[
A_*\ddot\eta+B\eta=0,\qquad \omega^2=\frac B{A_*}.
\]
$A_*$ 由所选坐标决定：直线位移对应等效质量，角度对应等效转动惯量。由角度改用弧长时，动能与势能的系数须同时变换，所得物理频率保持不变。
\begin{figure}[H]\centering
\begin{tikzpicture}
\draw[axis](-2.5,0)--(2.7,0)node[right]{$\eta$};\draw[axis](0,0)--(0,2.7)node[above]{$U-U_*$};
\draw[curve,domain=-1.8:1.8,samples=90]plot(\x,{.5*(\x)^2+.08*(\x)^4});\draw[guide,domain=-1.8:1.8,samples=80]plot(\x,{.5*(\x)^2});\node[right]at(1.6,.7){二次近似};
\end{tikzpicture}
\caption{势阱在平衡附近的二次近似。振幅增大后，高阶项可改变周期。}\label{sup4:fig:d7-potential}
\end{figure}
图\ref{sup4:fig:d7-potential}中，二次系数为正时存在非零的小振动频率。若二次项消失，不能把零频率当作稳定性结论，须保留高阶势能；若存在整体平移自由度，则零频率也可能对应匀速漂移。


\subsection{质量矩阵与固有振型}
将二次能量写为分量和式，利用 $M,K$ 的对称性，
\[
\frac{\partial T_2}{\partial\dot\eta_i}
=\frac12\sum_j(M_{ij}+M_{ji})\dot\eta_j=\sum_jM_{ij}\dot\eta_j,
\qquad
\frac{\partial U_2}{\partial\eta_i}=\sum_jK_{ij}\eta_j.
\]
质量矩阵在本阶取平衡点处常量，将动量对时间求导，拉格朗日方程给出 $\sum_jM_{ij}\ddot\eta_j+\sum_jK_{ij}\eta_j=0$，写成矩阵即
\begin{equation}
M\ddot{\vct\eta}+K\vct\eta=0.
\label{sup4:eq:d7-matrix}
\end{equation}
例如 $T_2$ 中若有 $c\dot\eta_1\dot\eta_2$，则 $M_{12}=M_{21}=c$，不是 $2c$；若势能含 $\kappa(\eta_2-\eta_1)^2/2$，展开后 $K_{12}=K_{21}=-\kappa$。

取所有坐标具有同一时间因子的解 $\vct\eta=\vct u\cos(\omega t+\phi)$，得到广义特征方程
\begin{equation}
(K-\omega^2M)\vct u=0,\qquad
\det(K-\omega^2M)=0.
\label{sup4:eq:d7-eigen}
\end{equation}
非零矢量 $\vct u$ 给出振型：各分量之比决定同一模态中各物体的振幅比例和相位关系。整体倍乘可并入模态振幅，故振型本身没有固定长度。
\begin{figure}[H]\centering
\begin{tikzpicture}
\draw[curve](0,0)--(7,0);\draw[curve](0,0)--(0,1.1);\draw[curve](7,0)--(7,1.1);\draw[spring](0,.5)--(2,.5);\draw[curve](2,.2)rectangle(2.7,.8);\draw[spring](2.7,.5)--(4.3,.5);\draw[curve](4.3,.2)rectangle(5,.8);\draw[spring](5,.5)--(7,.5);
\node at(2.35,.5){$m$};\node at(4.65,.5){$m$};\node[above]at(1,.65){$k$};\node[above]at(3.5,.65){$\kappa$};\node[above]at(6,.65){$k$};\draw[axis](2.1,-.5)--(3,-.5)node[right]{$x_1$};\draw[axis](4.4,-.5)--(5.3,-.5)node[right]{$x_2$};
\end{tikzpicture}
\caption{双振子以各自平衡位置为位移原点，两个正方向均向右。}\label{sup4:fig:d7-pair}
\end{figure}
图\ref{sup4:fig:d7-pair}的方程为
\[
m\ddot x_1+(k+\kappa)x_1-\kappa x_2=0,\qquad
m\ddot x_2+(k+\kappa)x_2-\kappa x_1=0.
\]
令 $x_+=x_1+x_2$、$x_-=x_1-x_2$。相加与相减分别给出
\[
m\ddot x_++kx_+=0,\qquad
m\ddot x_-+(k+2\kappa)x_-=0.
\]
所以同相振型 $(1,1)$ 和反相振型 $(1,-1)$ 的频率分别为 $\sqrt{k/m}$、$\sqrt{(k+2\kappa)/m}$。同相运动不拉伸中间弹簧，因此该模态的频率与 $\kappa$ 无关。


\subsection{例题与模型变式}

\begin{example}[圆环上的四个等质量球]\label{ex:hw4-1}

四个质量为$m$的小球可在固定光滑圆环上滑动，相邻两球用沿圆环布置的相同轻弹簧连接，劲度系数均为$k$。平衡时四球等距。求沿环方向小振动的全部简正模式。

\end{example}


\begin{figure}[H]\centering
\begin{tikzpicture}[>=Stealth,line cap=round,line join=round]
\draw(0,0)circle(1.5);\foreach \a/\i in {90/1,0/2,270/3,180/4}{\fill(\a:1.5)circle(3pt)node[shift={(\a:.35)}]{$\i$};}\foreach \a in {0,90,180,270}{\draw[spring,domain=\a+10:\a+80,samples=25]plot({1.5*cos(\x)},{1.5*sin(\x)});}\node at(0,0){$m,k$};
\end{tikzpicture}
\caption{由四个相同弹簧连接的环上质点}\label{fig:hw4-1-q}
\end{figure}

\begin{solution}
取图\ref{fig:hw4-1-s}中各球沿同一绕环正向的弧长位移$q_i$，下标按模4循环。弹簧伸长量为$q_{i+1}-q_i$，因此
\[
 m\ddot q_i=k(q_{i+1}+q_{i-1}-2q_i).
\]

\begin{figure}[H]\centering
\begin{tikzpicture}[>=Stealth,line cap=round,line join=round]
\node at(1.7,1.5){各$q_i$以顺时针为正，箭头表示模态位移};\begin{scope}[shift={(0,0)}]\draw(0,0)circle(.8);\fill(0.000000,0.800000)circle(2pt)node[shift={(90:.22)}]{$ 1 $};\draw[axis](0.000000,0.800000)--(0.500000,0.800000);\fill(0.800000,0.000000)circle(2pt)node[shift={(0:.22)}]{$ 2 $};\draw[axis](0.800000,0.000000)--(0.800000,-0.500000);\fill(-0.000000,-0.800000)circle(2pt)node[shift={(270:.22)}]{$ 3 $};\draw[axis](-0.000000,-0.800000)--(-0.500000,-0.800000);\fill(-0.800000,0.000000)circle(2pt)node[shift={(180:.22)}]{$ 4 $};\draw[axis](-0.800000,0.000000)--(-0.800000,0.500000);\node at(0,-1.55){$\omega^2=0$};\end{scope}\begin{scope}[shift={(3.5,0)}]\draw(0,0)circle(.8);\fill(0.000000,0.800000)circle(2pt)node[shift={(90:.22)}]{$ 1 $};\draw[axis](0.000000,0.800000)--(0.500000,0.800000);\fill(0.800000,0.000000)circle(2pt)node[shift={(0:.22)}]{$ 2 $};\fill(-0.000000,-0.800000)circle(2pt)node[shift={(270:.22)}]{$ 3 $};\draw[axis](-0.000000,-0.800000)--(0.500000,-0.800000);\fill(-0.800000,0.000000)circle(2pt)node[shift={(180:.22)}]{$ 4 $};\node at(0,-1.55){$\omega^2=2k/m$};\end{scope}\begin{scope}[shift={(0,-3)}]\draw(0,0)circle(.8);\fill(0.000000,0.800000)circle(2pt)node[shift={(90:.22)}]{$ 1 $};\fill(0.800000,0.000000)circle(2pt)node[shift={(0:.22)}]{$ 2 $};\draw[axis](0.800000,0.000000)--(0.800000,-0.500000);\fill(-0.000000,-0.800000)circle(2pt)node[shift={(270:.22)}]{$ 3 $};\fill(-0.800000,0.000000)circle(2pt)node[shift={(180:.22)}]{$ 4 $};\draw[axis](-0.800000,0.000000)--(-0.800000,-0.500000);\node at(0,-1.55){$\omega^2=2k/m$};\end{scope}\begin{scope}[shift={(3.5,-3)}]\draw(0,0)circle(.8);\fill(0.000000,0.800000)circle(2pt)node[shift={(90:.22)}]{$ 1 $};\draw[axis](0.000000,0.800000)--(0.500000,0.800000);\fill(0.800000,0.000000)circle(2pt)node[shift={(0:.22)}]{$ 2 $};\draw[axis](0.800000,0.000000)--(0.800000,0.500000);\fill(-0.000000,-0.800000)circle(2pt)node[shift={(270:.22)}]{$ 3 $};\draw[axis](-0.000000,-0.800000)--(-0.500000,-0.800000);\fill(-0.800000,0.000000)circle(2pt)node[shift={(180:.22)}]{$ 4 $};\draw[axis](-0.800000,0.000000)--(-0.800000,-0.500000);\node at(0,-1.55){$\omega^2=4k/m$};\end{scope}
\end{tikzpicture}
\caption{四个独立简正模及各球的相对位移方向}\label{fig:hw4-1-s}
\end{figure}

代入四个独立向量，可得：
\[
 \begin{array}{c|c}
 \vct a&\omega^2\\\hline
 (1,1,1,1)&0\\
 (1,0,-1,0)&2k/m\\
 (0,1,0,-1)&2k/m\\
 (1,-1,1,-1)&4k/m
 \end{array}
\]
第一种使各弹簧伸长均不变，是整体绕环的零频运动；其时间系数可为$A+Bt$。中间两个模式退化，其任意独立线性组合仍属于同一频率。最后一个模式中相邻球反向，回复作用最大。

四个向量构成完整基底，给定初始条件后可逐项投影。两个非零频率中有一个二重简并，因此仍须保留四个独立模式。


四个位移组成列向量后，质量矩阵为$m\vct 1$，刚度矩阵为
\[
 \vct K=k\begin{pmatrix}
 2&-1&0&-1\\-1&2&-1&0\\0&-1&2&-1\\-1&0&-1&2
 \end{pmatrix}.
\]
令$a_j=\exp[i\kappa(j-1)]$，周期条件要求$\kappa=s\pi/2$，$s=0,1,2,3$。代入方程即得$m\omega^2=2k(1-\cos\kappa)$，从而得到表中的零频、二重简并频率和最高频率。复指数的实部、虚部分别对应两个实的简并模式。

若初位移、初速度分别为$\vct q_0,\vct w_0$，表中任一非零频模式$\vct a_s$的系数为
\[
 Q_s(t)=\frac{\vct a_s\cdot\vct q_0}{\vct a_s\cdot\vct a_s}\cos\omega_st
 +\frac{\vct a_s\cdot\vct w_0}{\omega_s\vct a_s\cdot\vct a_s}\sin\omega_st.
\]
零频模式则有$Q_0(t)=\tfrac14\sum q_{i0}+\tfrac14\sum w_{i0}\,t$。把$Q_s(t)\vct a_s$相加，便得到满足任意小初态的完整运动。图\ref{fig:hw4-1-s}中的箭头均按绕环的同一正向判断，不能直接比较箭头在纸面上的左右。

\end{solution}



\begin{example}[两定滑轮与三等质量物块]\label{ex:hw4-2}

图\ref{fig:hw4-2-q}中两个光滑定滑轮等高，三物块质量均为$m$，中间物块由两条不可伸长轻绳悬挂，两端各挂一个物块。平衡时两段斜绳长均为$l$，与竖直夹角$60^\circ$。求平面内的两个小振动模式。

\end{example}


\begin{figure}[H]\centering
\begin{tikzpicture}[>=Stealth,line cap=round,line join=round]
\coordinate(A)at(-2,0);\coordinate(B)at(2,0);\coordinate(C)at(0,-1.15);\draw(A)circle(.15);\draw(B)circle(.15);\draw[curve](-2.15,-1.8)--(-2.15,0)--(A)--(C)--(B)--(2.15,0)--(2.15,-1.8);\draw[fill=white](-2.4,-2.2)rectangle(-1.9,-1.8);\node at(-2.15,-2){$m$};\draw[fill=white](1.9,-2.2)rectangle(2.4,-1.8);\node at(2.15,-2){$m$};\draw[fill=white](-0.25,-1.3499999999999999)rectangle(0.25,-0.95);\node at(0,-1.15){$m$};\node[above]at(A){$A$};\node[above]at(B){$B$};\node at(-.75,-.3){$l$};\node at(.75,-.3){$l$};
\end{tikzpicture}
\caption{两定滑轮悬挂三个等质量物块}\label{fig:hw4-2-q}
\end{figure}

\begin{solution}
取中间物块相对平衡位置的位移为$(x,y)$，$y$向下。两斜段长度为
\[
 r_{1,2}=\sqrt{(\pm a+x)^2+(h+y)^2},\quad
 a=\sqrt3l/2,\ h=l/2.
\]
图\ref{fig:hw4-2-s}两端物块向下位移$z_i=-(r_i-l)$。一阶关系给出
\[
 z_1=-\frac{\sqrt3}{2}x-\frac12y,\quad
 z_2=\frac{\sqrt3}{2}x-\frac12y.
\]

\begin{figure}[H]\centering
\begin{tikzpicture}[>=Stealth,line cap=round,line join=round]
\coordinate(A)at(-2,0);\coordinate(B)at(2,0);\coordinate(C)at(0,-1.15);\draw(A)circle(.15);\draw(B)circle(.15);\draw[curve](-2.15,-1.8)--(-2.15,0)--(A)--(C)--(B)--(2.15,0)--(2.15,-1.8);\draw[fill=white](-2.4,-2.2)rectangle(-1.9,-1.8);\node at(-2.15,-2){$m$};\draw[fill=white](1.9,-2.2)rectangle(2.4,-1.8);\node at(2.15,-2){$m$};\draw[fill=white](-0.25,-1.3499999999999999)rectangle(0.25,-0.95);\node[left]at(-.4,-1.45){$m$};\node[above]at(A){$A$};\node[above]at(B){$B$};\node at(-.75,-.3){$l$};\node at(.75,-.3){$l$};\draw[axis](0,-1.15)--(1,-1.15)node[above]{$x$};\draw[axis](0,-1.15)--(0,-2.2)node[left]{$y$};\node at(0,.65){$a=\sqrt3l/2,\quad h=l/2$};
\end{tikzpicture}
\caption{中央物块的两个独立小位移}\label{fig:hw4-2-s}
\end{figure}

把根号写为$r_{1,2}=l\sqrt{1+\varepsilon_{1,2}}$，其中
\[
\varepsilon_{1,2}=\frac{\pm\sqrt3x+y}{l}+\frac{x^2+y^2}{l^2}.
\]
用$\sqrt{1+\varepsilon}=1+\varepsilon/2-\varepsilon^2/8+O(\varepsilon^3)$展开；二次项中的$\varepsilon^2$只需保留$(\pm\sqrt3x+y)^2/l^2$。合并同次项后，绳长为
\[
 r_{1,2}=l\pm\frac{\sqrt3}{2}x+\frac12y
 +\frac{x^2+3y^2\mp2\sqrt3xy}{8l}+O(3).
\]
动能只需一阶位移关系；势能的一次项互相抵消后，二阶绳长才产生回复作用。

由此$T=\frac m2(\frac52\dot x^2+\frac32\dot y^2)$。势能须把绳长展开到二阶，得到
\[
 U=-mg(y+z_1+z_2)
 =\frac{mg}{4l}x^2+\frac{3mg}{4l}y^2+\text{常量}.
\]
分别对二次动能、势能求导，有
\[
\frac{\dd}{\dd t}\frac{\partial T}{\partial\dot x}
=\frac{5m}{2}\ddot x,\quad
\frac{\partial U}{\partial x}=\frac{mg}{2l}x;
\qquad
\frac{\dd}{\dd t}\frac{\partial T}{\partial\dot y}
=\frac{3m}{2}\ddot y,\quad
\frac{\partial U}{\partial y}=\frac{3mg}{2l}y.
\]
因此两方程为$\ddot x+gx/(5l)=0$与$\ddot y+gy/l=0$。各自的特征根为$\pm\mathrm i\sqrt{g/(5l)}$和$\pm\mathrm i\sqrt{g/l}$，故
\[
 \boxed{\omega_x=\sqrt{g/(5l)},\qquad\omega_y=\sqrt{g/l}}.
\]
水平模式的两端物块反向升降；竖直模式的两端物块同向，且中间物块位移与它们的位移方向相反、大小为两倍。若将绳长只展开到一阶，势能中这些回复项就会遗漏。


对水平纯模式取$y=0$，两端速度为$\dot z_{1,2}=\mp\sqrt3\dot x/2$；对竖直纯模式取$x=0$，两端速度均为$-\dot y/2$。这两组位移比例给出了图\ref{fig:hw4-2-s}中三个物块在每个模式下的实际运动。

\end{solution}



\begin{example}[半密度矩形长木条的浮动]\label{ex:hw4-4}

均匀矩形长木条密度为水的一半，横截面水平宽度$a$、竖直高度$b$，长度$L$。木条浮在水面。(1)求水平平衡稳定的条件；(2)求上下平动与绕长轴摆动的周期。忽略流体附加质量及波浪辐射，按静水压力计算回复作用。

\end{example}


\begin{figure}[H]\centering
\begin{tikzpicture}[>=Stealth,line cap=round,line join=round]
\draw[thick](-2,0)--(2.5,0);\fill[pattern=north east lines](-2,-1.6)rectangle(2.5,-1.5);\draw[curve](-1,-.7)rectangle(1,.7);\draw[guide](-1,0)--(1,0);\fill(0,0)circle(1pt)node[above left]{$C$};\node at(2,0.35){水面};\draw[<->](-1,-1.05)--(1,-1.05)node[midway,below]{$a$};\draw[<->](1.4,-.7)--(1.4,.7)node[midway,right]{$b$};
\end{tikzpicture}
\caption{半密度木条的平衡吃水}\label{fig:hw4-4-q}
\end{figure}

\begin{solution}
平衡时吃水$b/2$，质量$m=\rho_wabL/2$。图\ref{fig:hw4-4-s}中下沉位移$z$使排水体积增加$aLz$，故
\[
 m\ddot z=-\rho_wgaLz,\qquad
 \boxed{T_z=2\pi\sqrt{b/(2g)}}.
\]

\begin{figure}[H]\centering
\begin{tikzpicture}[>=Stealth,line cap=round,line join=round]
\draw[thick](-2.1,0)--(2.1,0);\node[right]at(2.1,0){水面};
\begin{scope}\clip(-2,-1.6)rectangle(2,0);\begin{scope}[rotate=15]\fill[gray!18](-1.2,-.7)rectangle(1.2,.7);\end{scope}\end{scope}
\begin{scope}[rotate=15]\draw[curve](-1.2,-.7)rectangle(1.2,.7);\draw[guide](0,-.85)--(0,1.7);\end{scope}
\fill(0,0)circle(1.5pt)node[above right,fill=white,inner sep=1pt]{$C$};
\fill(-.09,-.35)circle(1.2pt)node[left,fill=white,inner sep=1pt]{$B$};\fill(-.09,.336)circle(1.2pt)node[left,fill=white,inner sep=1pt]{$M$};
\draw[axis](-.09,-.35)--(-.09,1.25)node[right]{$F_b$};\draw[axis](0,0)--(0,-1.15)node[right]{$mg$};
\draw[guide](0,0)--(0,1.65);\draw(0,1.25)arc(90:105:1.25);\node at(-.5,1.45){$\theta$};
\node at(0,-1.85){$x_B\simeq-CM\theta$};
\begin{scope}[shift={(5.1,0)}]\draw[thick](-1.2,0)--(1.2,0);\draw[curve](-.8,-1.2)rectangle(.8,.2);
\draw[guide](-.8,-.7)--(.8,-.7);\draw[<->](1.2,-.7)--(1.2,-1.2)node[midway,right]{$z$};
\node at(0,.9){下沉分量};\node at(0,-1.6){$\Delta V=aLz$};\end{scope}
\end{tikzpicture}
\caption{倾斜时浮心偏移产生回复力矩，下沉时排水体积增加}\label{fig:hw4-4-s}
\end{figure}

小倾角 $\theta$ 下，浮心初始在质心下方 $b/4$。水线面的面积二阶矩为 $L\int_{-a/2}^{a/2}x^2\dd x=L[x^3/3]_{-a/2}^{a/2}=La^3/12$，除以排水体积 $abL/2$，得到稳心半径 $a^2/(6b)$。因此稳心高为
\[
 CM=\frac{a^2}{6b}-\frac b4.
\]
浮心的横向偏移也可由排水区域积分求出。若截面逆时针倾斜$\theta>0$，在随木条转动的坐标中，水面为$y=-\theta x+O(\theta^2)$，浮心的横坐标为
\[
 x_B^{\rm body}=\frac{-\theta\int_{-a/2}^{a/2}x^2\dd x}{ab/2}
 =-\frac{a^2}{6b}\theta.
\]
转换回地面坐标时还要计入浮心原在质心下方$b/4$，故$x_B=-CM\theta$。图\ref{fig:hw4-4-s}中的浮力作用线通过偏移后的$B$，其关于$C$的力矩为$-mgCM\theta$。$I_C=m(a^2+b^2)/12$，故
\[
 \boxed{\frac ba<\sqrt{\frac23},\qquad
 T_\theta=2\pi\sqrt{\frac{b(a^2+b^2)}{g(2a^2-3b^2)}}}.
\]
这里的开区间保证非零的线性回复频率。等号处可继续保留高阶项：只要水线仍切过两侧边，水面在随体坐标中为 $y=-x\tan\theta$，排水截面面积仍为 $ab/2$。横向一阶矩只需将前式的 $\theta$ 换为 $\tan\theta$；竖向一阶矩为
\[
\int_{-a/2}^{a/2}\int_{-b/2}^{-x\tan\theta}y\dd y\dd x
=\frac12\int_{-a/2}^{a/2}\left(x^2\tan^2\theta-\frac{b^2}4\right)\dd x
=\frac{a^3}{24}\tan^2\theta-\frac{ab^2}8.
\]
这里的内层原函数为 $y^2/2$，外层使用上文已算出的 $x^2$ 积分。除以排水截面面积，浮心在随体坐标中的位置为
\[
 x_B^{\rm body}=-\frac{a^2}{6b}\tan\theta,\qquad
 y_B^{\rm body}=-\frac b4+\frac{a^2}{12b}\tan^2\theta.
\]
将其转回地面系，浮力矩为
\[
 \tau_C=-mg\sin\theta\left[CM+\frac{a^2}{12b}\tan^2\theta\right].
\]
因此 $CM=0$ 时，$\tau_C=-mga^2\theta^3/(12b)+O(\theta^5)$。由 $\dd U/\dd\theta=-\tau_C$，积分得
\[
U(\theta)-U(0)=\frac{mga^2}{12b}\int_0^\theta\xi^3\dd\xi+O(\theta^6)
=\frac{mga^2}{12b}[\xi^4/4]_0^\theta+O(\theta^6).
\]
势能的首个转动项为 $mga^2\theta^4/(48b)>0$。在题设静水模型中，完整的局部稳定条件为 $b/a\le\sqrt{2/3}$；等号处没有有限的线性周期，不能套用上面的 $T_\theta$。超过该比值时水平构形不稳定。

对称性使上下与转动在一阶近似中解耦。实际流体的附加惯性会改变周期，上述周期取决于题设采用的静水模型。


\end{solution}



\begin{example}[两均匀杆组成的复摆]\label{ex:hw4-6}

两根均匀细杆长度均为$l$、质量均为$m$，第一杆上端铰接于固定点$O$，下端与第二杆上端在$A$处铰接。两杆在同一竖直面内绕向下平衡位置作小振动，求简正模式。

\end{example}


\begin{figure}[H]\centering
\begin{tikzpicture}[>=Stealth,line cap=round,line join=round]
\draw[thick](-.5,0)--(.7,0);\coordinate(O)at(0,0);\coordinate(A)at(.4,-1.7);\coordinate(B)at(1.1,-3.3);\draw[curve](O)--(A)--(B);\fill(A)circle(2pt)node[right]{$A$};\node[above]at(O){$O$};\node at(-.3,-.8){$\theta_1$};\node at(.2,-2.5){$\theta_2$};\draw[guide](O)--(0,-1.7);\draw[guide](A)--(.4,-3.3);
\end{tikzpicture}
\caption{两均匀杆铰接成的复摆}\label{fig:hw4-6-q}
\end{figure}

\begin{solution}
图\ref{fig:hw4-6-s}中$\theta_1,\theta_2$均相对竖直向下量取，为两杆的绝对倾角。两杆质心位置分别为$(l\sin\theta_1/2,-l\cos\theta_1/2)$与$(l\sin\theta_1+l\sin\theta_2/2,-l\cos\theta_1-l\cos\theta_2/2)$。

\begin{figure}[H]\centering
\begin{tikzpicture}[>=Stealth,line cap=round,line join=round]
\draw[thick](-.5,0)--(.7,0);\coordinate(O)at(0,0);\coordinate(A)at(.4,-1.7);\coordinate(B)at(1.1,-3.3);\draw[curve](O)--(A)--(B);\fill(A)circle(2pt)node[right]{$A$};\node[above]at(O){$O$};\node at(-.3,-.8){$\theta_1$};\node at(.2,-2.5){$\theta_2$};\draw[guide](O)--(0,-1.7);\draw[guide](A)--(.4,-3.3);\draw[axis](.2,-.85)--(1,-.85)node[above]{$v_{C1,x}^{(1)}$};\draw[axis](.75,-2.5)--(1.65,-2.5)node[above]{$v_{C2,x}^{(1)}$};\node at(3,-1.5){$I_C=ml^2/12$};
\end{tikzpicture}
\caption{质心速度取到位移一阶，水平分量进入二阶动能}\label{fig:hw4-6-s}
\end{figure}

第一杆关于固定上端的动能为$ml^2\dot\theta_1^2/6$。第二杆质心的水平速度在一阶近似中为$l\dot\theta_1+l\dot\theta_2/2$，再加上自身转动动能，得
\[
 T_2=\frac{ml^2}{2}\dot\theta_1^2+
 \frac{ml^2}{2}\dot\theta_1\dot\theta_2+
 \frac{ml^2}{6}\dot\theta_2^2.
\]
交叉项来自第二杆质心同时受到上杆运动和自身转动的影响。

将质心平动及各杆绕质心的转动相加，保留二次项：
\[
 T=\frac{ml^2}{2}\left(\frac43\dot\theta_1^2+\dot\theta_1\dot\theta_2+\frac13\dot\theta_2^2\right),
 \quad U=\frac{mgl}{2}\left(\frac32\theta_1^2+\frac12\theta_2^2\right).
\]
由于二次动能不显含角坐标，拉格朗日方程中的惯性项为
\[
\frac{\dd}{\dd t}\frac{\partial T}{\partial\dot\theta_1}
=ml^2\left(\frac43\ddot\theta_1+\frac12\ddot\theta_2\right),\qquad
\frac{\dd}{\dd t}\frac{\partial T}{\partial\dot\theta_2}
=ml^2\left(\frac12\ddot\theta_1+\frac13\ddot\theta_2\right).
\]
势能导数分别为$3mgl\theta_1/2$、$mgl\theta_2/2$。令$\theta_j=a_j\cos\omega t$，将$\ddot\theta_j=-\omega^2\theta_j$代入两个方程。再置$\lambda=\omega^2l/g$，非零振幅要求
\[
 \det\begin{pmatrix}3/2-4\lambda/3&-\lambda/2\\-\lambda/2&1/2-\lambda/3\end{pmatrix}=0
 \quad\Longrightarrow\quad7\lambda^2-42\lambda+27=0.
\]
故
\[
 \boxed{\omega_\pm^2=\left(3\pm\frac6{\sqrt7}\right)\frac gl,
 \qquad\frac{a_2}{a_1}=\frac{3\lambda_\pm}{3-2\lambda_\pm}}.
\]
低频根给出同向摆动，高频根给出反向摆动。若把$\theta_2$误当成两杆夹角，质心速度和动能交叉项都会改变；方程与图中的角度定义须始终一致。


展开行列式，得
\[
 \frac34-\frac76\lambda+\left(\frac49-\frac14\right)\lambda^2=0.
\]
乘以36即为$7\lambda^2-42\lambda+27=0$。求振幅比时采用第二行$-(\lambda/2)a_1+(1/2-\lambda/3)a_2=0$，便得$3\lambda/(3-2\lambda)$。低频根$\lambda_-\simeq0.7322$的振幅比约为$1.4305$，高频根$\lambda_+\simeq5.2678$的振幅比约为$-2.0972$。图\ref{fig:hw4-6-s}中的两杆倾角都以固定竖直线为零点，故正比值表示同向，负比值表示反向。

若给定初态，也应使用质量矩阵$ml^2\left(\begin{smallmatrix}4/3&1/2\\1/2&1/3\end{smallmatrix}\right)$投影。第二杆随第一杆运动，使质量矩阵具有非零的非对角元；仅按两杆质量相等作普通投影并不成立。

\end{solution}



\begin{example}[摆杆、弹簧与悬重的耦合]\label{ex:hw4-7}

长$l$、质量$3m$的均匀杆上端固定铰接，平衡时竖直。杆下端左侧通过劲度系数$k$的水平弹簧连接墙，右侧通过同样弹簧及绕定滑轮的轻绳连接质量$m$的悬重。弹簧预伸长使题设构形平衡，求小振动简正模式。

\end{example}


\begin{figure}[H]\centering
\begin{tikzpicture}[>=Stealth,line cap=round,line join=round]
\draw[thick](-.5,0)--(.5,0);\draw[curve](0,0)--(0,-2);\draw[spring](-2,-2)--(0,-2);\draw[spring](0,-2)--(1.75,-2);\draw(2,-2.15)circle(.15);\draw(1.75,-2)--(2,-2)arc(90:0:.15)--(2.15,-3.3);\draw[thick](-2,-2.5)--(-2,-1.5);\draw[fill=white](1.9,-3.7)rectangle(2.4,-3.3);\node at(2.15,-3.5){$m$};\node at(.6,-1){$3m,l$};\node at(-1,-1.5){$k$};\node at(1,-1.5){$k$};
\end{tikzpicture}
\caption{弹簧连接的均匀杆与悬挂物块}\label{fig:hw4-7-q}
\end{figure}

\begin{solution}
令杆向右小转角为$\theta$、悬重相对平衡位置向下位移为$x=l\varphi$。图\ref{fig:hw4-7-s}右弹簧的伸长增量为$x-l\theta$，左弹簧为$l\theta$。

\begin{figure}[H]\centering
\begin{tikzpicture}[>=Stealth,line cap=round,line join=round]
\draw[thick](-.5,0)--(.5,0);\draw[curve](0,0)--(0,-2);\draw[spring](-2,-2)--(0,-2);\draw[spring](0,-2)--(1.75,-2);\draw(2,-2.15)circle(.15);\draw(1.75,-2)--(2,-2)arc(90:0:.15)--(2.15,-3.3);\draw[thick](-2,-2.5)--(-2,-1.5);\draw[fill=white](1.9,-3.7)rectangle(2.4,-3.3);\node at(2.15,-3.5){$m$};\node at(.6,-1){$3m,l$};\node at(-1,-1.5){$k$};\node at(1,-1.5){$k$};\draw[axis](0,-2.35)--(.8,-2.35)node[below]{$l\theta$};\draw[axis](2.8,-3)--(2.8,-3.8)node[right]{$x$};
\end{tikzpicture}
\caption{杆端水平位移与物块竖直位移}\label{fig:hw4-7-s}
\end{figure}

杆的惯量为$(3m)l^2/3=ml^2$。以平衡位置为展开点，一次项由预拉力和重力抵消，二次式为
\[
 T=\frac12ml^2(\dot\theta^2+\dot\varphi^2),\quad
 U=\frac{3mgl}{4}\theta^2+\frac{kl^2}{2}[\theta^2+(\varphi-\theta)^2].
\]
对势能求导，
\[
\frac{\partial U}{\partial\theta}
=\frac{3mgl}{2}\theta+kl^2(2\theta-\varphi),\qquad
\frac{\partial U}{\partial\varphi}=kl^2(\varphi-\theta).
\]
动能分别给出$ml^2\ddot\theta$、$ml^2\ddot\varphi$，故设$a=3g/(2l)$、$b=k/m$后，两方程为
\[
\ddot\theta+(a+2b)\theta-b\varphi=0,\qquad
\ddot\varphi+b\varphi-b\theta=0.
\]
取$\theta=A\cos\omega t$、$\varphi=B\cos\omega t$，代入得
\[
 \begin{pmatrix}a+2b-\omega^2&-b\\-b&b-\omega^2\end{pmatrix}
 \binom{A}{B}=0.
\]
令$\lambda=\omega^2$，行列式逐项展开为
\[
(a+2b-\lambda)(b-\lambda)-b^2
=\lambda^2-(a+3b)\lambda+ab+b^2=0.
\]
二次方程的判别式为$(a+3b)^2-4(ab+b^2)=(a+b)^2+4b^2$。用求根公式，并从矩阵第二行$-bA+(b-\lambda)B=0$求振幅比，得到
\[
 \boxed{\omega_\pm^2=\frac{a+3b\pm\sqrt{(a+b)^2+4b^2}}2,
 \qquad B/A=\frac{b}{b-\omega_\pm^2}}.
\]
两个频率平方之和为$a+3b>0$，乘积为$ab+b^2>0$，且判别式为$(a+b)^2+4b^2$，故$k>0$时两根均为正。低频模式中$B/A>0$，杆端向右时悬重向下；高频模式中两者反向。

这里$B/A$是$\varphi/\theta$的振幅比，实际悬重位移与角振幅之比为$lB/A$。只有再给定$k/m=g/l$之类的附加关系，才能进一步简化频率；这里保留两个独立参数。


由特征方程$(a+2b-\omega^2)(b-\omega^2)-b^2=0$，可知低频根小于$b$而高频根大于$b$，因而上述振幅比的符号与图示一致。求弱弹簧极限时，令$\varepsilon=b/a\ll1$。频率公式的根式为
\[
\sqrt{(a+b)^2+4b^2}=a\sqrt{1+2\varepsilon+5\varepsilon^2}.
\]
在$\sqrt{1+z}=1+z/2-z^2/8+O(z^3)$中代入$z=2\varepsilon+5\varepsilon^2$，二次系数为$5/2-4/8=2$，故根式等于$a+b+2b^2/a+O(b^3/a^2)$。分别代入两根，得到
\[
 \omega_-^2=b-\frac{b^2}{a}+O(b^3/a^2),\qquad
 \omega_+^2=a+2b+\frac{b^2}{a}+O(b^3/a^2).
\]
前者以悬重运动为主，杆的振幅较小；后者趋近单独均匀杆的物理摆频率$\sqrt{3g/(2l)}$。这里极限$b\to0^+$只描述二次振动方程的参数变化；若实际装置不能维持题设预拉力，则须另行确定新的平衡构形。

\end{solution}



\begin{example}[正方形顶点电荷的平面稳定性]\label{ex:hw4-8}

边长$L$的正方形两条对角线端点上，分别放置两对固定正电荷$Q$与$\alpha Q$，$\alpha>0$。质量$m$、电荷$q>0$的粒子静止在中心，运动限制在正方形平面内。求稳定平衡的$\alpha$范围及两个方向的小振动周期。

\end{example}


\begin{figure}[H]\centering
\begin{tikzpicture}[>=Stealth,line cap=round,line join=round]
\draw[curve](-1.6,0)--(0,1.6)--(1.6,0)--(0,-1.6)--cycle;\foreach \x/\y/\lab in {-1.6/0/Q,1.6/0/Q,0/1.6/\alpha Q,0/-1.6/\alpha Q}{\fill(\x,\y)circle(2pt)node[shift={({\x*.18},{\y*.18})}]{$\lab$};}\fill(0,0)circle(2pt)node[below right]{$q,m$};
\end{tikzpicture}
\caption{正方形四顶点的交替电荷}\label{fig:hw4-8-q}
\end{figure}

\begin{solution}
图\ref{fig:hw4-8-s}沿两条对角线取$x,y$轴，两对电荷坐标分别为$(\pm a,0)$和$(0,\pm a)$，其中$a=L/\sqrt2$。为展开每一对电荷的势，令$\vct e$沿该对电荷的轴线、$r^2=x^2+y^2$。对每一项使用$(1+z)^{-1/2}=1-z/2+3z^2/8+O(z^3)$，有
\[
\frac1{|a\vct e\mp\vct r|}
=\frac1a\left[1\pm\frac{\vct e\cdot\vct r}{a}
+\frac{3(\vct e\cdot\vct r)^2-r^2}{2a^2}\right]+O(r^3).
\]
正负两项相加消去奇次项，再分别取$\vct e=\vct e_x,\vct e_y$并乘相应电荷量，得到
\[
 U=\frac{2k_eqQ(1+\alpha)}a
 +\frac{k_eqQ}{a^3}[(2-\alpha)x^2+(2\alpha-1)y^2]+O(|\vct r|^4).
\]

\begin{figure}[H]\centering
\begin{tikzpicture}[>=Stealth,line cap=round,line join=round]
\draw[curve](-1.6,0)--(0,1.6)--(1.6,0)--(0,-1.6)--cycle;\foreach \x/\y in {-1.6/0,1.6/0,0/1.6,0/-1.6}{\fill(\x,\y)circle(2pt);}\node at(-1.95,.2){$Q$};\node at(1.95,.2){$Q$};\node[right]at(.16,1.95){$\alpha Q$};\node at(0,-2){$\alpha Q$};\fill(0,0)circle(2pt)node[below right]{$q,m$};\draw[axis](0,0)--(2.5,0)node[above]{$x$};\draw[axis](0,0)--(0,2.4)node[right]{$y$};\node at(3,-1){$a=L/\sqrt2$};
\end{tikzpicture}
\caption{平面内两条对角方向取为坐标轴}\label{fig:hw4-8-s}
\end{figure}

一次项因成对对称抵消，$xy$项也为零。两个线性回复系数都为正的条件是$1/2<\alpha<2$。这个开区间给出两个非零的小振动频率；判断稳定平衡的完整范围还须检查端点。
对二次势能分别求导，线性运动方程为
\[
m\ddot x=-\frac{\partial U_2}{\partial x}
=-\frac{2k_eqQ(2-\alpha)}{a^3}x,\qquad
m\ddot y=-\frac{\partial U_2}{\partial y}
=-\frac{2k_eqQ(2\alpha-1)}{a^3}y.
\]
各自除以$m$，括号内位移的系数就是$\omega_x^2,\omega_y^2$。在此开区间内，由$T=2\pi/\omega$得到
\[
 \boxed{T_x=2\pi\sqrt{\frac{ma^3}{2k_eqQ(2-\alpha)}},\quad
 T_y=2\pi\sqrt{\frac{ma^3}{2k_eqQ(2\alpha-1)}}}.
\]
沿$x$轴，四电荷的精确势能可先合并为
\[
U(x,0)=k_eqQ\left[\frac{2a}{a^2-x^2}
+\frac{2\alpha}{\sqrt{a^2+x^2}}\right].
\]
第一项用$(1-z)^{-1}=1+z+z^2+O(z^3)$，第二项用上述二项式，均取$z=x^2/a^2$。减去中心值后，
\[
U(x,0)-U(0,0)=\frac{k_eqQ}{a}
\left[(2-\alpha)\frac{x^2}{a^2}
+\left(2+\frac{3\alpha}{4}\right)\frac{x^4}{a^4}\right]+O(x^6).
\]
在$\alpha=2$时，二次项消失，留下
\[
 U(x,0)-U(0,0)=\frac{7k_eqQ}{2a^5}x^4+O(x^6)>0.
\]
沿$y$轴则有$U(0,y)=k_eqQ[2\alpha a/(a^2-y^2)+2/\sqrt{a^2+y^2}]$；按相同的两种二项式展开，二次项系数为$2\alpha-1$、四次项系数为$2\alpha+3/4$。代入$\alpha=1/2$，得到
\[
 U(0,y)-U(0,0)=\frac{7k_eqQ}{4a^5}y^4+O(y^6)>0.
\]
还须证明任意平面扰动都会增加势能，而不只检查两条坐标轴。以$\alpha=2$为例，记$A_y=3k_eqQ/a^3>0$、$B_x=7k_eqQ/(2a^5)>0$。势能关于$x,y$分别为偶函数，故在原点邻域可写成
\[
U-U(0,0)=A_y y^2+B_x x^4+R(x,y),
\]
其中未写出的四次项只有$x^2y^2,y^4$，更高项从六次开始。因此存在有限常数$C_1,C_2,C_3$，使
\[
|R|\le C_1x^2y^2+C_2y^4+C_3(x^2+y^2)^3.
\]
在$|x|,|y|\le\delta$内，分别有$x^2y^2\le\delta^2y^2$、$y^4\le\delta^2y^2$，并且
\[
(x^2+y^2)^3=x^6+3x^4y^2+3x^2y^4+y^6
\le4\delta^2x^4+4\delta^4y^2.
\]
取$\delta$足够小，可使$|R|\le B_xx^4/2+A_yy^2/2$，于是
\[
U-U(0,0)\ge\frac{A_y}{2}y^2+\frac{B_x}{2}x^4>0
\quad\text{（非零扰动）}.
\]
这证明$\alpha=2$仍为严格局部极小。$\alpha=1/2$时交换$x,y$，并将势能乘一个正常数，即回到同一论证。因此完整的平面稳定范围为
\[
 \boxed{\tfrac12\le\alpha\le2}.
\]
端点的退化方向由三次回复力控制，振动周期随振幅趋于零而发散；上述有限周期公式只用于开区间。本题明确限制在平面内；离开平面时四个斥力势都随$|z|$增加而降低，中心并非三维稳定平衡。


\end{solution}



\begin{example}[圆弧槽内的纯滚动振动]\label{deep4:problem-example-1}\label{deep4:ex:d7-e1}
半径为 $R$ 的固定圆弧槽开口向上，半径为 $a<R$、质量为 $m$ 的圆形刚体在槽内无滑动滚动，绕自身轴的转动惯量为 $I=\beta ma^2$。以最低点为 $\theta=0$，刚体中心与槽圆心的连线相对竖直向下方向的转角为 $\theta$。求小振动频率。若从 $\theta=\Theta$ 静止释放，$0<\Theta<\pi/2$，求维持全过程纯滚动所需的最小静摩擦系数。
\end{example}
\begin{figure}[H]\centering
\begin{tikzpicture}
\draw[curve](-2.954,-.521)arc(-170:-10:3);\fill(0,0)circle(1.5pt)node[above]{$O$};\draw[guide](0,0)--(0,-1.25);\coordinate(C)at({2.5*sin(30)},{-2.5*cos(30)});\coordinate(S)at({3*sin(30)},{-3*cos(30)});\draw[curve](C)circle(.5);\fill(C)circle(1.7pt);\draw[guide](C)--(.1,-1.9)node[left]{$C$};\fill(S)circle(1.5pt)node[below right]{$S$};\draw[guide](0,0)--(S);\node[left]at(.65,-1.15){$\ell$};\draw[<->](C)--++(.5,0)node[midway,above]{$a$};\draw(0,-.65)arc(-90:-60:.65);\node at(.5,-.7){$\theta$};\node at(-1.9,-.6){固定圆弧槽};
\end{tikzpicture}
\caption{中心C沿半径 $\ell=R-a$ 的圆运动，接触点S位于OC的延长线上。}\label{deep4:fig:d7-e1-q}
\end{figure}
\begin{solution}
图\ref{deep4:fig:d7-e1-q}中，中心速率为 $\ell|\dot\theta|$，其中 $\ell=R-a$。接触点相对固定槽瞬时静止，因此刚体自转角速度大小为 $\ell|\dot\theta|/a$。总动能与势能为
\[
T=\frac12m\ell^2\dot\theta^2+\frac12I\frac{\ell^2\dot\theta^2}{a^2}
=\frac12m(1+\beta)\ell^2\dot\theta^2,
\qquad U=mg\ell(1-\cos\theta).
\]
为求转折点处也适用的运动方程，按图\ref{deep4:fig:d7-e1-s}取增大 $\theta$ 的切向为正，用滚动约束消去自转角加速度，平动与转动方程成为
\[
m\ell\ddot\theta=-mg\sin\theta+f,\qquad
\beta m\ell\ddot\theta=-f.
\]
静摩擦同时改变平动和自转，虽不耗散能量，仍会影响振动频率。
两式相加得到
\[
(1+\beta)\ell\ddot\theta+g\sin\theta=0.
\]
在线性近似中 $\sin\theta\simeq\theta$，故
\[
\boxed{\omega=\sqrt{\frac{g}{(1+\beta)(R-a)}}.}
\]
等效惯量包含平动与自转两部分，所得频率是同轨道上光滑滑动质点频率的 $1/\sqrt{1+\beta}$。
\begin{figure}[H]\centering
\begin{tikzpicture}
\fill(0,0)circle(2pt)node[below left]{$C$};\draw[guide](-1,1.732)--(1,-1.732);\draw[axis](0,0)--(-.9,1.559)node[above]{$N$};\draw[axis](0,0)--(0,-1.8)node[below]{$mg$};\draw[axis](0,0)--(1.2,.693)node[right]{$f$};\node[align=left]at(4.1,.7){$\theta>0$ 时\\摩擦沿槽向上\\正切向取增大 $\theta$ 的方向};
\end{tikzpicture}
\caption{法向力指向槽圆心，摩擦提供滚动所需的自转角加速度。}\label{deep4:fig:d7-e1-s}
\end{figure}
按图\ref{deep4:fig:d7-e1-s}取沿槽向上为切向正方向。平动方程为 $m\ell\ddot\theta=-mg\sin\theta+f$，代入已求出的 $\ddot\theta$，得
\[
f=\frac{\beta}{1+\beta}mg\sin\theta.
\]
中心的向心加速度给出 $N-mg\cos\theta=m\ell\dot\theta^2$。从转折点释放后的能量关系为
\[
\ell\dot\theta^2=\frac{2g}{1+\beta}(\cos\theta-\cos\Theta),
\quad
N=\frac{mg}{1+\beta}[(\beta+3)\cos\theta-2\cos\Theta].
\]
由于 $|\theta|\le\Theta<\pi/2$，$N>0$，接触始终维持。摩擦比值为
\[
\frac{|f|}{N}=\frac{\beta|\sin\theta|}{(\beta+3)\cos\theta-2\cos\Theta}.
\]
该比值与习题\ref{deep1:exercise:d1-p5}的摩擦比是同一函数，只需取 $\kappa=\beta$、$\theta_0=\Theta$。使用那里已完整展开的商导数，在 $0\le\theta\le\Theta$ 上，
\[
\frac{\dd}{\dd\theta}\frac{|f|}{N}
=\frac{\beta[(\beta+3)-2\cos\Theta\cos\theta]}
{[(\beta+3)\cos\theta-2\cos\Theta]^2}>0.
\]
其余半程关于最低点对称，故最大值在两个转折点取得。因此
\[
\boxed{\mu_{s,\min}=\frac{\beta}{1+\beta}\tan\Theta.}
\]
摩擦条件在给定的 $0<\Theta<\pi/2$ 范围内未使用小角近似，频率公式则要求 $\Theta\ll1$。对实心球，$\beta=2/5$，小振动频率为 $\sqrt{5g/[7(R-a)]}$。
\end{solution}



\section{振型分解、简正坐标与初始条件}

频率不能单独确定一次实际运动，初位移和初速度还要分别分解到振型中。不同质量的系统采用质量加权的正交关系，再由模态运动恢复各物体的位移。



\subsection{振型正交与初始条件分解}
若 $K\vct u_r=\omega_r^2M\vct u_r$，由矩阵对称性可得
\[
\vct u_s^{\mathsf T}K\vct u_r
=\omega_r^2\vct u_s^{\mathsf T}M\vct u_r
=\omega_s^2\vct u_s^{\mathsf T}M\vct u_r.
\]
当两频率不同，
\begin{equation}
\vct u_s^{\mathsf T}M\vct u_r=0.
\label{sup4:eq:d7-orthogonal}
\end{equation}
正交关系的权重是质量矩阵。质量不相等或动能含交叉项时，投影必须使用这个权重。选取归一化振型 $\vct e_r$，使 $\vct e_s^{\mathsf T}M\vct e_r=\delta_{sr}$；对简并频率，也可在相应子空间内按此内积选取正交基。

写成 $\vct\eta=\sum_r\vct e_r Q_r$，则
\[
\ddot Q_r+\omega_r^2Q_r=0,\qquad
Q_r(0)=\vct e_r^{\mathsf T}M\vct\eta(0),\quad
\dot Q_r(0)=\vct e_r^{\mathsf T}M\dot{\vct\eta}(0).
\]
对 $\omega_r>0$，特征方程 $\lambda^2+\omega_r^2=0$ 有根 $\pm\mathrm i\omega_r$。实通解为 $Q_r=A_r\cos\omega_rt+B_r\sin\omega_rt$，由初态得 $A_r=Q_r(0)$、$B_r=\dot Q_r(0)/\omega_r$，即
\[
Q_r(t)=Q_r(0)\cos\omega_rt+\frac{\dot Q_r(0)}{\omega_r}\sin\omega_rt.
\]
若未归一化，投影式须再除以 $\vct u_r^{\mathsf T}M\vct u_r$。每个模态能量独立守恒，因为 $\dd E_r/\dd t=\dot Q_r(\ddot Q_r+\omega_r^2Q_r)=0$：
\[
E_2-U_* =\frac12\sum_r(\dot Q_r^2+\omega_r^2Q_r^2).
\]
\begin{figure}[H]\centering
\begin{tikzpicture}
\begin{scope}
\draw[axis](-1.7,0)--(1.9,0)node[right]{$x_1$};\draw[axis](0,-1.4)--(0,1.6)node[above]{$x_2$};\draw[axis](0,0)--(1.2,1.2)node[above right]{$\vct u_+$};\draw[axis](0,0)--(1.2,-1.2)node[below right]{$\vct u_-$};\node at(0,-2){等质量的正交振型};
\end{scope}
\begin{scope}[xshift=5.5cm]
\draw[axis](-1.8,0)--(2,0)node[right]{$Q_+$};\draw[axis](0,-1.5)--(0,1.7)node[above]{$Q_-$};\draw[curve](0,0)ellipse(1.3 and .85);\node at(0,-2){模态坐标中的能量椭圆};
\end{scope}
\end{tikzpicture}
\caption{左图以等质量双振子示意振型方向；右图为给定能量、速度为零截面上的椭圆。}\label{sup4:fig:d7-modes}
\end{figure}
图\ref{sup4:fig:d7-modes}的模态分解也可用于冲量后的运动。广义冲量为 $\vct J_q$ 时，$M\Delta\dot{\vct\eta}=\vct J_q$。$\vct J_q$ 须由作用点的虚位移求出；对角坐标，它还包含相应的力臂。


\subsection{简正坐标的构造与退化情形}
\label{sec:complete4}
质量矩阵正定时，可写成$M=S^{\mathsf T}S$，令$\vct z=S\vct\eta$，则
\[
 T_2=\frac12\dot{\vct z}^{\mathsf T}\dot{\vct z},\qquad
 U_2-U_* =\frac12\vct z^{\mathsf T}
 (S^{-\mathsf T}KS^{-1})\vct z.
\]
括号内的实对称矩阵可经正交变换对角化，因而$\omega_j^2$为实数，振型可选成完备的质量正交基。坐标由长度改成角度时，质量矩阵和刚度矩阵须同时变换；单独比较某个矩阵的对角元没有坐标不变的意义。

\subsubsection{由运动和初态求模态能量}
若采用未归一化振型$\vct u_j$，记$M_j=\vct u_j^{\mathsf T}M\vct u_j$，并写$\vct\eta=\sum_j\vct u_jQ_j$。利用正交关系，得到
\[
 E_j=\frac12M_j(\dot Q_j^2+\omega_j^2Q_j^2),\quad
 Q_j(0)=\frac{\vct u_j^{\mathsf T}M\vct\eta(0)}{M_j}.
\]
初速度按相同权重投影。若系统原来静止，受到广义冲量$\vct J_q$，则$\dot Q_j(0^+)=\vct u_j^{\mathsf T}\vct J_q/M_j$。若受冲量的位置是某一振型的节点，该模态的初速度投影为零。

\subsection{例题与模型变式}

\begin{example}[弹簧耦合的双单摆]\label{ex:hw4-3}

两个质量$m$的小球分别由长$l$的轻杆悬挂在等高、相距$l$的两铰点下，球之间用原长$l$、劲度系数$k$的轻弹簧相连。初态为静止平衡，瞬间给$A$球水平速度$v_0$。求小振动范围内两球的运动。

\end{example}


\begin{figure}[H]\centering
\begin{tikzpicture}[>=Stealth,line cap=round,line join=round]
\draw[thick](-.6,0)--(2.6,0);\draw[curve](0,0)--(0,-2);\draw[curve](2,0)--(2,-2);\fill(0,-2)circle(2pt)node[below]{$A,m$};\fill(2,-2)circle(2pt)node[below]{$B,m$};\draw[spring](0,-2)--(2,-2);\node at(1,-1.6){$k,l$};\node at(-.3,-1){$l$};\draw[<->](0,.45)--(2,.45)node[midway,above]{$l$};
\end{tikzpicture}
\caption{等长双摆在平衡位置由弹簧连接}\label{fig:hw4-3-q}
\end{figure}

\begin{solution}
用水平小位移$x_1,x_2$，图\ref{fig:hw4-3-s}中杆角约为$x_i/l$，弹簧伸长为$x_2-x_1$。运动方程为
\[
 m\ddot x_1=-\frac{mg}{l}x_1-k(x_1-x_2),\quad
 m\ddot x_2=-\frac{mg}{l}x_2-k(x_2-x_1).
\]

\begin{figure}[H]\centering
\begin{tikzpicture}[>=Stealth,line cap=round,line join=round]
\draw[thick](-.6,0)--(2.6,0);\draw[curve](0,0)--(0,-2);\draw[curve](2,0)--(2,-2);\fill(0,-2)circle(2pt)node[below]{$A,m$};\fill(2,-2)circle(2pt)node[below]{$B,m$};\draw[spring](0,-2)--(2,-2);\node at(1,-1.6){$k,l$};\node at(-.3,-1){$l$};\draw[<->](0,.45)--(2,.45)node[midway,above]{$l$};\draw[axis](0,-2.4)--(.7,-2.4)node[below]{$x_1$};\draw[axis](2,-2.4)--(2.7,-2.4)node[below]{$x_2$};
\end{tikzpicture}
\caption{两个水平位移确定弹簧伸长量}\label{fig:hw4-3-s}
\end{figure}

相加、相减得$q_s=x_1+x_2$、$q_a=x_1-x_2$，满足
\[
 \ddot q_s+\omega_s^2q_s=0,\quad\omega_s^2=g/l,
 \qquad\ddot q_a+\omega_a^2q_a=0,\quad\omega_a^2=g/l+2k/m.
\]
两方程的特征根分别为$\pm\mathrm i\omega_s$、$\pm\mathrm i\omega_a$，通解写成
\[
q_s=A_s\cos\omega_st+B_s\sin\omega_st,\qquad
q_a=A_a\cos\omega_at+B_a\sin\omega_at.
\]
初态$q_s(0)=q_a(0)=0$给出$A_s=A_a=0$；对通解求导并代入$\dot q_s(0)=\dot q_a(0)=v_0$，得$\omega_sB_s=\omega_aB_a=v_0$。再由$x_1=(q_s+q_a)/2$、$x_2=(q_s-q_a)/2$还原，得到
\[
 \boxed{x_1=\frac{v_0}{2}\left(\frac{\sin\omega_st}{\omega_s}
 +\frac{\sin\omega_at}{\omega_a}\right),\quad
 x_2=\frac{v_0}{2}\left(\frac{\sin\omega_st}{\omega_s}
 -\frac{\sin\omega_at}{\omega_a}\right)}.
\]
由上式可直接验证$x_1(0)=x_2(0)=0$、$\dot x_1(0)=v_0$、$\dot x_2(0)=0$。一个足够保证小振幅的条件是
\[
 \frac{v_0}{2l}\left(\frac1{\omega_s}+\frac1{\omega_a}\right)\ll1.
\]
若不满足这一条件，球的水平位移不再等于$l\theta_i$，弹簧长度也不能只保留$x_2-x_1$。

同向模式不改变弹簧长度，反向模式则同时受到重力和弹簧回复。两频接近时可见缓慢的振幅交换，但精确解仍是上式的两个正弦叠加。


\end{solution}



\begin{example}[三弹簧与不等质量双振子]\label{ex:hw4-5}

两小球在光滑水平面上由三个相同弹簧依次与两墙连接，弹簧劲度系数为$k$。分别求$m_1=m_2=m$及$m_1=m,m_2=2m$时的简正频率与振幅比。

\end{example}


\begin{figure}[H]\centering
\begin{tikzpicture}[>=Stealth,line cap=round,line join=round]
\draw[fill=white](0.75,-0.2)rectangle(1.25,0.2);\node at(1,0){$m_1$};\draw[fill=white](2.75,-0.2)rectangle(3.25,0.2);\node at(3,0){$m_2$};\draw[spring](-.5,0)--(.75,0);\draw[spring](1.25,0)--(2.75,0);\draw[spring](3.25,0)--(4.5,0);\draw[thick](-.5,-.5)--(-.5,.5);\draw[thick](4.5,-.5)--(4.5,.5);\node at(2,.7){三个弹簧均为$k$};
\end{tikzpicture}
\caption{三弹簧连接的两个质量块}\label{fig:hw4-5-q}
\end{figure}

\begin{solution}
图\ref{fig:hw4-5-s}的$x_1,x_2$从各自平衡位置向右量起，回复力为$-2kx_1+kx_2$及$kx_1-2kx_2$。因此
\[
 \det\begin{pmatrix}2k-m_1\omega^2&-k\\-k&2k-m_2\omega^2\end{pmatrix}=0.
\]

\begin{figure}[H]\centering
\begin{tikzpicture}[>=Stealth,line cap=round,line join=round]
\draw[fill=white](0.75,-0.2)rectangle(1.25,0.2);\node at(1,0){$m_1$};\draw[fill=white](2.75,-0.2)rectangle(3.25,0.2);\node at(3,0){$m_2$};\draw[spring](-.5,0)--(.75,0);\draw[spring](1.25,0)--(2.75,0);\draw[spring](3.25,0)--(4.5,0);\draw[thick](-.5,-.5)--(-.5,.5);\draw[thick](4.5,-.5)--(4.5,.5);\node at(2,.7){三个弹簧均为$k$};\draw[axis](1,-.5)--(1.8,-.5)node[below]{$x_1$};\draw[axis](3,-.5)--(3.8,-.5)node[below]{$x_2$};
\begin{scope}[shift={(1,-1.8)}]
\fill(0,0)circle(2pt)node[below]{$m$};\fill(2,0)circle(2pt)node[below]{$2m$};
\draw[axis](0,0)--(.55,0);\draw[axis](2,0)--(2.7513,0);\node[left]at(-.7,0){低频};
\node at(1,-.6){$a_2/a_1=(1+\sqrt3)/2$};\end{scope}
\begin{scope}[shift={(1,-3.2)}]
\fill(0,0)circle(2pt)node[below]{$m$};\fill(2,0)circle(2pt)node[below]{$2m$};
\draw[axis](0,0)--(.55,0);\draw[axis](2,0)--(1.7987,0);\node[left]at(-.7,0){高频};
\node at(1,-.6){$a_2/a_1=(1-\sqrt3)/2$};\end{scope}
\end{tikzpicture}
\caption{不等质量系统中，两模式的位移方向和相对振幅}\label{fig:hw4-5-s}
\end{figure}

等质量时，两模式为$(1,1)$、$(1,-1)$，分别有$\boxed{\omega_s^2=k/m,\ \omega_a^2=3k/m}$。

不等质量时置$\lambda=m\omega^2/k$，行列式化为$2\lambda^2-6\lambda+3=0$，故
\[
 \boxed{\omega_\pm^2=\frac{3\pm\sqrt3}{2}\frac km,\qquad
 \frac{a_2}{a_1}=2-\lambda_\pm=\frac{1\mp\sqrt3}{2}}.
\]
取$\vct a_\pm=(1,(1\mp\sqrt3)/2)$。它们的普通内积不为零，但
\[
 \vct a_-^{\mathsf T}\begin{pmatrix}m&0\\0&2m\end{pmatrix}\vct a_+
 =m\left[1+2\frac{1+\sqrt3}{2}\frac{1-\sqrt3}{2}\right]=0.
\]
因此展开初态时应采用质量加权投影，才能得到独立模式的振幅。

低频模式中两球同向，高频模式中反向。此时简单的$x_1\pm x_2$不再自动解耦；可令$\vct x=c_-\vct a_-\cos(\omega_-t+\delta_-)+c_+\vct a_+\cos(\omega_+t+\delta_+)$，再用初态求系数。


设$\vct M=\operatorname{diag}(m,2m)$，给定初态为$\vct x_0,\vct v_0$。用质量内积将两式分别投影，得到
\[
 q_s(0)=\frac{\vct a_s^{\mathsf T}\vct M\vct x_0}{\vct a_s^{\mathsf T}\vct M\vct a_s},\qquad
 \dot q_s(0)=\frac{\vct a_s^{\mathsf T}\vct M\vct v_0}{\vct a_s^{\mathsf T}\vct M\vct a_s},
 \quad s=\pm.
\]
于是$\vct x(t)=\sum_s\vct a_s[q_s(0)\cos\omega_st+\dot q_s(0)\sin\omega_st/\omega_s]$。例如只把第一球拉开$d$再静止释放，$q_s(0)=d/[1+2r_s^2]$，其中$r_s=(1\mp\sqrt3)/2$。两个系数之和为$d$，而$r_-q_-(0)+r_+q_+(0)=0$，恰好还原两球的初位移。

\end{solution}



\begin{example}[弹簧小车上的单摆]\label{sup4:problem-example-2}\label{sup4:ex:d7-e2}
质量为 $M$ 的小车在光滑水平面上运动，通过劲度系数为 $k$ 的弹簧连接固定墙。车上悬挂长为 $\ell$、摆球质量为 $m$ 的单摆。车的平衡位移记为 $x$，摆角 $\theta$ 从竖直向下方向向右量起。求两个小振动频率和振型，讨论 $k\to\infty$ 的极限。
\end{example}
\begin{figure}[H]\centering
\begin{tikzpicture}
\draw[curve](0,0)--(6,0);\draw[curve](0,0)--(0,1.5);\draw[spring](0,.6)--(2,.6)node[midway,above]{$k$};\draw[curve](2,.2)rectangle(3.8,1);\node at(2.9,.55){$M$};\draw[curve](2.9,1)--(2.9,2.8);\fill(2.9,2.8)circle(1.5pt)node[above]{$P$};\draw[curve](2.9,2.8)--(3.7,1.3)node[midway,right]{$\ell$};\fill(3.7,1.3)circle(2.7pt)node[right]{$m$};\draw[guide](2.9,2.8)--(2.9,1.3);\draw(2.9,2.15)arc(-90:-62:.65);\node at(3.15,1.9){$\theta$};\draw[axis](2.3,-.45)--(3.4,-.45)node[right]{$x$};
\end{tikzpicture}
\caption{P随小车平移。摆球的地面水平位移一阶为 $x+\ell\theta$。}\label{sup4:fig:d7-e2-q}
\end{figure}
\begin{solution}
图\ref{sup4:fig:d7-e2-q}中，摆球地面坐标为 $(x+\ell\sin\theta,-\ell\cos\theta)$。因此二次动能和势能为
\[
T_2=\frac12(M+m)\dot x^2+m\ell\dot x\dot\theta
+\frac12m\ell^2\dot\theta^2,
\qquad U_2=\frac12kx^2+\frac12mg\ell\theta^2.
\]
摆球的精确速率平方为 $\dot x^2+2\ell\cos\theta\,\dot x\dot\theta+\ell^2\dot\theta^2$。其中交叉项说明：小车速度与摆动速度同向时，摆球动能增加，反向时则减小。保留二次项后，由动能、势能逐项列出
\[
\begin{aligned}
(M+m)\ddot x+m\ell\ddot\theta+kx&=0,\\
m\ell\ddot x+m\ell^2\ddot\theta+mg\ell\theta&=0.
\end{aligned}
\]
第一式是整个车摆系统的水平动量方程，第二式描述加速悬点下的摆动。对一支模态取 $x=X\cos\omega t$、$\theta=\Theta\cos\omega t$，将二阶导数换成 $-\omega^2$ 倍的位移，即得下面的矩阵方程。
对应矩阵为
\[
M_q=\begin{pmatrix}M+m&m\ell\\m\ell&m\ell^2\end{pmatrix},
\qquad K_q=\begin{pmatrix}k&0\\0&mg\ell\end{pmatrix}.
\]
记 $\lambda=\omega^2$，特征行列式除以 $m\ell^2$ 后得到
\[
M\lambda^2-\left[k+\frac{(M+m)g}{\ell}\right]\lambda
+\frac{kg}{\ell}=0.
\]
行列式展开的具体形式为
\[
[k-(M+m)\lambda](g-\ell\lambda)-m\ell\lambda^2=0.
\]
其判别式可写成 $[k-(M+m)g/\ell]^2+4mkg/\ell>0$，两根之和、之积又分别为正，故确有两个不同的正频率平方。
故两频率平方为
\[
\boxed{\omega_{\pm}^2=
\frac{k+(M+m)g/\ell\ \pm\sqrt{[k+(M+m)g/\ell]^2-4Mkg/\ell}}{2M}.}
\]
将 $\lambda=g/\ell$ 代入上面的二次多项式，值为 $-m(g/\ell)^2<0$；而两个根均为正，所以 $\omega_-^2<g/\ell<\omega_+^2$。第二行模态方程给出
\[
\frac{\Theta}{X}=\frac{\omega^2}{\ell(g/\ell-\omega^2)}.
\]
\begin{figure}[H]\centering
\begin{tikzpicture}
\begin{scope}
\draw[curve](0,0)rectangle(1.7,.65);\draw[curve](.15,.65)--(.15,2.85)--(.85,2.85);\fill(.85,2.85)circle(1.5pt)node[above right]{$P$};
\draw[guide](.85,2.85)--(.85,1.45);\draw[curve](.85,2.85)--(1.3,1.6);\fill(1.3,1.6)circle(2pt);
\draw[axis](.2,3.45)--(1.3,3.45)node[right]{$X$};\draw[axis](.85,1.3)--(1.3,1.3)node[midway,below]{$\ell\Theta$};\node at(.9,-.5){低频：$X\Theta>0$};
\end{scope}
\begin{scope}[xshift=4.7cm]
\draw[curve](0,0)rectangle(1.7,.65);\draw[curve](.15,.65)--(.15,2.85)--(.85,2.85);\fill(.85,2.85)circle(1.5pt)node[above right]{$P$};
\draw[guide](.85,2.85)--(.85,1.45);\draw[curve](.85,2.85)--(.4,1.6);\fill(.4,1.6)circle(2pt);
\draw[axis](.2,3.45)--(1.3,3.45)node[right]{$X$};\draw[axis](.85,1.3)--(.4,1.3)node[midway,below]{$\ell\Theta$};\node at(.9,-.5){高频：$X\Theta<0$};
\end{scope}
\end{tikzpicture}
\caption{两种模态的位移方向。$X$ 为小车位移，$\ell\Theta$ 为摆球相对悬点P的水平位移；箭头长度仅作示意。}\label{sup4:fig:d7-e2-s}
\end{figure}
图\ref{sup4:fig:d7-e2-s}中，低频模态的小车位移与摆角同号，高频模态中两者反号。摆球在地面系中的水平位移为 $x+\ell\theta$，其中 $\ell\theta$ 只表示相对悬点的位移。

求$k\to\infty$的极限时，不直接计算低频根中两个大数的差。将高频根的根式除以$k$，可得
\[
\frac{\omega_+^2}{k}
=\frac{1+(M+m)g/(k\ell)+
\sqrt{[1+(M+m)g/(k\ell)]^2-4Mg/(k\ell)}}{2M}
\longrightarrow\frac1M.
\]
再由两根之积$\omega_-^2\omega_+^2=kg/(M\ell)$，得到
\[
\omega_-^2=\frac{kg}{M\ell\omega_+^2}\longrightarrow\frac g\ell.
\]
对低频模态保持有限的摆角振幅，振幅比公式给出$X\to0$，恢复固定支点单摆。高频模态则有$\ell\Theta/X=\omega_+^2/(g/\ell-\omega_+^2)\to-1$，摆球的地面水平位移近乎抵消。弹簧在这一快速模态中主要驱动小车自身，所以$\omega_+^2\sim k/M$。
\end{solution}



\begin{example}[等长等质量双摆]\label{sup4:problem-example-6}\label{sup4:ex:d7-e6}
双摆两杆均轻且长为 $\ell$，两质点质量均为 $m$，在同一竖直平面运动。两摆角 $\theta_1,\theta_2$ 均从竖直向下方向量起。求小振动频率及振型，并求初始 $\theta_1=\Theta,\theta_2=0$、两角速度均为零时的运动。
\end{example}
\begin{figure}[H]\centering
\begin{tikzpicture}
\fill(0,0)circle(1.6pt)node[above]{$O$};\coordinate(A)at(.8,-1.8);\coordinate(B)at(.3,-3.7);\draw[curve](0,0)--(A)--(B);\fill(A)circle(2.5pt)node[right]{$m$};\fill(B)circle(2.5pt)node[left]{$m$};\draw[guide](0,0)--(0,-1.5);\draw[guide](A)--(.8,-3.3);\draw(0,-.7)arc(-90:-66:.7);\node at(.22,-1.25){$\theta_1$};\draw(.8,-2.5)arc(-90:-105:.7);\node at(.66,-3.1){$\theta_2$};\node[right]at(.5,-1){$\ell$};\node[left]at(.4,-2.7){$\ell$};\node[align=left]at(4,-1.8){两角均为绝对摆角\\第二角不从第一杆量起};
\end{tikzpicture}
\caption{第二杆图示向左偏，故 $\theta_2<0$。两个正角均取向右偏转。}\label{sup4:fig:d7-e6-q}
\end{figure}
\begin{solution}
按图\ref{sup4:fig:d7-e6-q}，两质点水平位移的一阶式分别为 $\ell\theta_1$、$\ell(\theta_1+\theta_2)$。竖直速度只从二阶开始，因此
\[
T_2=\frac12m\ell^2(2\dot\theta_1^2+2\dot\theta_1\dot\theta_2+\dot\theta_2^2),
\qquad U_2=\frac12mg\ell(2\theta_1^2+\theta_2^2).
\]
竖直位移对势能仍须保留二次项。上质点上升 $\ell\theta_1^2/2$，下质点上升 $\ell(\theta_1^2+\theta_2^2)/2$，相加便得到 $U_2$。由二次能量列出的两个角坐标方程为
\[
2\ddot\theta_1+\ddot\theta_2+\frac{2g}{\ell}\theta_1=0,\qquad
\ddot\theta_1+\ddot\theta_2+\frac g\ell\theta_2=0.
\]
图\ref{sup4:fig:d7-e6-q}中的第二角是第二杆相对竖直线的绝对角；若改成两杆夹角，坐标关系和矩阵都会改变。
两个矩阵为
\[
M=m\ell^2\begin{pmatrix}2&1\\1&1\end{pmatrix},\qquad
K=mg\ell\begin{pmatrix}2&0\\0&1\end{pmatrix}.
\]
令 $\eta=\omega^2\ell/g$，特征方程为
\[
\det\begin{pmatrix}2-2\eta&-\eta\\-\eta&1-\eta\end{pmatrix}
=\eta^2-4\eta+2=0.
\]
故
\[
\boxed{\omega_-^2=(2-\sqrt2)g/\ell,\qquad
\omega_+^2=(2+\sqrt2)g/\ell.}
\]
由第二行 $-\eta\Theta_1+(1-\eta)\Theta_2=0$，得到低频振型 $(1,\sqrt2)$、高频振型 $(1,-\sqrt2)$。
\begin{figure}[H]\centering
\begin{tikzpicture}
\begin{scope}
\fill(0,0)circle(1.5pt);\draw[curve](0,0)--(.45,-1.5)--(1.07,-2.95);\fill(.45,-1.5)circle(2pt);\fill(1.07,-2.95)circle(2pt);\node at(.5,.55){低频};\node at(.5,-3.5){$\Theta_2=\sqrt2\Theta_1$};
\end{scope}
\begin{scope}[xshift=4.8cm]
\fill(0,0)circle(1.5pt);\draw[curve](0,0)--(.45,-1.5)--(-.17,-2.95);\fill(.45,-1.5)circle(2pt);\fill(-.17,-2.95)circle(2pt);\node at(.3,.55){高频};\node at(.3,-3.5){$\Theta_2=-\sqrt2\Theta_1$};
\end{scope}
\end{tikzpicture}
\caption{两模态中，两杆绝对摆角同号或反号。构形只示小角比例。}\label{sup4:fig:d7-e6-s}
\end{figure}
图\ref{sup4:fig:d7-e6-s}的比例针对两杆摆角。第二质点的水平位移为 $\ell(\theta_1+\theta_2)$，故两质点带符号的水平振幅比为 $(\Theta_1+\Theta_2)/\Theta_1$。

若两个余弦模态的系数为 $c_-,c_+$，初态分别要求 $c_-+c_+=\Theta$、$\sqrt2c_--\sqrt2c_+=0$，所以两者均为 $\Theta/2$。两初速度为零，又使两个正弦系数同时为零。
初始构形分解为 $(\Theta,0)=\frac\Theta2(1,\sqrt2)+\frac\Theta2(1,-\sqrt2)$，因此
\[
\boxed{\theta_1=\frac\Theta2(\cos\omega_-t+\cos\omega_+t),\qquad
\theta_2=\frac\Theta{\sqrt2}(\cos\omega_-t-\cos\omega_+t).}
\]
本题势能矩阵已对角，耦合来自动能交叉项。第一杆偏转还同时抬高两个质点，这使 $\theta_1^2$ 的势能系数为单个摆球的两倍。
\end{solution}



\begin{example}[三振子链的节点与冲量选择]\label{deep4:problem-example-9}\label{deep4:ex:d7-e9}
三个相同质量 $m$ 的滑块依次由四根相同弹簧 $k$ 与两端固定墙连接，在光滑直线上运动。求全部小振动模态。系统初始位于平衡，对中间块施加向右冲量 $J$，求此后的运动及各模态能量。
\end{example}
\begin{figure}[H]\centering
\begin{tikzpicture}
\draw[curve](0,0)--(9,0);\draw[curve](0,0)--(0,1.2);\draw[curve](9,0)--(9,1.2);\foreach\x in{1.6,4.1,6.6}{\draw[curve](\x,.2)rectangle++(.8,.7);\node at(\x+.4,.55){$m$};}\draw[spring](0,.55)--(1.6,.55);\draw[spring](2.4,.55)--(4.1,.55);\draw[spring](4.9,.55)--(6.6,.55);\draw[spring](7.4,.55)--(9,.55);\foreach\x in{.8,3.25,5.75,8.2}{\node[above]at(\x,.7){$k$};}\draw[axis](4.5,1.2)--(5.8,1.2)node[right]{$J$};\foreach\x/\i in{2/1,4.5/2,7/3}{\node[below]at(\x,0){$x_\i$};}
\end{tikzpicture}
\caption{中间块受冲量；左右两块在初态关于中点对称。}\label{deep4:fig:d7-e9-q}
\end{figure}
\begin{solution}
图\ref{deep4:fig:d7-e9-q}的刚度矩阵为
\[
K=k\begin{pmatrix}2&-1&0\\-1&2&-1\\0&-1&2\end{pmatrix},\qquad M=mI_3.
\]
矩阵中的对角元2来自每个物块连接的两根弹簧，邻接元 $-1$ 来自中间弹簧形变量之差。若直接求特征行列式，令 $\eta=m\omega^2/k$，有
\[
\det(K/k-\eta I_3)=(2-\eta)[(2-\eta)^2-2].
\]
利用图\ref{deep4:fig:d7-e9-s}的镜面对称性，可以先分出反对称模态，再在对称子空间求另外两根，避免展开三次多项式。
反对称振型可先由镜面对称性找出：$(1,0,-1)$，中间块不动，频率平方为 $2k/m$。其余模态左右同向，置振型为 $(1,b,1)$，外块和中间块方程给出
\[
2-b=\eta,\qquad 2b-2=\eta b,\qquad \eta=m\omega^2/k.
\]
消元得 $b^2=2$。按频率升序，
\[
\begin{array}{c|c|c}
&\omega_j^2&\vct v_j\ \text{（普通单位长度）}\\\hline
1&(2-\sqrt2)k/m&(1,\sqrt2,1)/2\\
2&2k/m&(1,0,-1)/\sqrt2\\
3&(2+\sqrt2)k/m&(1,-\sqrt2,1)/2
\end{array}
\]
\begin{figure}[H]\centering
\begin{tikzpicture}
\foreach\y/\j in{2/1,0/2,-2/3}{\node[left]at(0,\y){模态\j};\foreach\x in{1,3.5,6}{\fill(\x,\y)circle(2pt);}}
\draw[axis](1,2)--(1.8,2);\draw[axis](3.5,2)--(4.631,2);\draw[axis](6,2)--(6.8,2);
\draw[axis](1,0)--(1.8,0);\node[above]at(3.5,0){节点};\draw[axis](6,0)--(5.2,0);
\draw[axis](1,-2)--(1.8,-2);\draw[axis](3.5,-2)--(2.369,-2);\draw[axis](6,-2)--(6.8,-2);
\end{tikzpicture}
\caption{中间块在模态2中为节点，打在此处的冲量不激起该模态。}\label{deep4:fig:d7-e9-s}
\end{figure}
冲量后的速度为 $(0,J/m,0)$。令 $\vct x=\sum\vct v_j a_j(t)$，因振型已按普通长度归一化且三质量相同，模态等效质量均为m。投影得到
\[
\dot a_1(0)=J/(\sqrt2m),\quad \dot a_2(0)=0,\quad
\dot a_3(0)=-J/(\sqrt2m),\qquad a_j(0)=0.
\]
冲量先确定模态初速度。对普通单位振型 $\vct v_j$，投影给出 $m\dot a_j(0)=\vct v_j^{\mathsf T}(0,J,0)^{\mathsf T}$；再解 $\ddot a_j+\omega_j^2a_j=0$，得到 $a_j(t)=\dot a_j(0)\sin\omega_jt/\omega_j$。位移振幅中的频率分母由这一步产生。
所以
\[
\boxed{\begin{aligned}
x_1=x_3&=\frac{J}{2\sqrt2m}\left(\frac{\sin\omega_1t}{\omega_1}-\frac{\sin\omega_3t}{\omega_3}\right),\\
x_2&=\frac{J}{2m}\left(\frac{\sin\omega_1t}{\omega_1}+\frac{\sin\omega_3t}{\omega_3}\right).
\end{aligned}}
\]
图\ref{deep4:fig:d7-e9-s}的模态2具有中间节点，初始位移、模态速度均为零，此后始终不出现。各模态能量为 $E_1=E_3=J^2/(4m)$、$E_2=0$，总和为中间块初始动能 $J^2/(2m)$。
\end{solution}



\section{拍振、失谐与外力响应}

近邻频率的叠加产生慢包络，周期外力则按其施加方式激发不同振型。自由拍振和强迫响应分别讨论，以区分能量交换、频率避交及反共振。



\subsection{拍振、周期外力与反共振}
两个相近频率叠加时，利用
\[
\cos\omega_1t+\cos\omega_2t
=2\cos\frac{\omega_2-\omega_1}{2}t\,
\cos\frac{\omega_2+\omega_1}{2}t,
\]
可把运动分为快速振动和缓慢包络。只有 $|\omega_2-\omega_1|\ll(\omega_1+\omega_2)/2$ 时，才有清楚分离的快、慢时间尺度。包络过零的时刻，位移为零，速度未必也为零；中间耦合弹簧还可能储存能量。

若外力为 $\vct f\cos\nu t$，周期特解 $\vct X\cos\nu t$ 满足
\begin{equation}
(K-\nu^2M)\vct X=\vct f.
\label{sup4:eq:d7-forced}
\end{equation}
当 $\nu$ 接近固有频率且外力对该模态的投影非零，会出现共振。若某个坐标的响应分子为零，而矩阵仍可逆，该坐标的周期响应可为零，称为反共振。无阻尼系统从任意初态出发还包含不衰减的自由振动；上述周期特解须与初始条件一起使用。
\begin{figure}[H]\centering
\begin{tikzpicture}
\draw[axis](0,0)--(6.1,0)node[right]{$\nu^2$};\draw[axis](0,-1.6)--(0,2)node[above]{有符号响应};
\draw[guide](1.5,-1.5)--(1.5,1.8);\draw[guide](4.5,-1.5)--(4.5,1.8);\node[below]at(1.5,-1.5){$\omega_1^2$};\node[below]at(4.5,-1.5){$\omega_2^2$};
\draw[curve](2,-1.3)..controls(2.6,-.5)and(3.3,.5)..(4,1.3);\fill(3,0)circle(1.8pt)node[below right]{反共振};\node at(3,2){两共振频率之间的一种情形};
\end{tikzpicture}
\caption{反共振是某个观测坐标的周期响应过零，其他坐标仍可运动。曲线仅示意零点与极点的次序。}\label{sup4:fig:d7-antiresonance}
\end{figure}
图\ref{sup4:fig:d7-antiresonance}中的极点还须结合外力和观测位置判断。若外力恰与某模态正交，就不能激起该模态；若测点位于该振型的节点，该频率也可能不出现在测点响应中。


\subsection{例题与模型变式}

\begin{example}[对称双振子与拍振包络]\label{sup4:problem-example-3}\label{sup4:ex:d7-e3}
两个质量均为 $m$ 的滑块在光滑直线上运动，左右外侧分别由弹簧 $k$ 连到固定墙，中间由弹簧 $\kappa$ 相连，三弹簧在平衡态均无形变。初始 $x_1=A,x_2=0$，速度均为零。求运动与模态能量，讨论弱耦合时的拍振，并判断包络过零是否代表第一块已经静止。
\end{example}
\begin{figure}[H]\centering
\begin{tikzpicture}
\draw[curve](0,0)--(7.5,0);\draw[curve](0,0)--(0,1.2);\draw[curve](7.5,0)--(7.5,1.2);\draw[spring](0,.55)--(2.3,.55);\draw[curve](2.3,.2)rectangle(3,.9);\draw[spring](3,.55)--(4.5,.55);\draw[curve](4.5,.2)rectangle(5.2,.9);\draw[spring](5.2,.55)--(7.5,.55);\node at(2.65,.55){$m$};\node at(4.85,.55){$m$};\node[above]at(1.1,.65){$k$};\node[above]at(3.75,.65){$\kappa$};\node[above]at(6.3,.65){$k$};\node at(2.65,-.55){$x_1(0)=A$};\node at(5.2,-.55){$x_2(0)=0$};
\end{tikzpicture}
\caption{两个位移均向右为正，图示初始位移条件。}\label{sup4:fig:d7-e3-q}
\end{figure}
\begin{solution}
图\ref{sup4:fig:d7-e3-q}中间弹簧的形变量为 $x_2-x_1$。它对第一块的力为 $\kappa(x_2-x_1)$，对第二块为相反数；外侧弹簧各提供 $-kx_i$。所以
\[
\begin{aligned}
m\ddot x_1&=-kx_1+\kappa(x_2-x_1),\\
m\ddot x_2&=-kx_2-\kappa(x_2-x_1).
\end{aligned}
\]
两式相加时耦合力抵消，相减时耦合力加倍，这就是选取同相、反相坐标的依据。
图\ref{sup4:fig:d7-e3-q}的同相与反相坐标取为 $s=(x_1+x_2)/2$、$d=(x_1-x_2)/2$。由两物块方程相加、相减，
\[
\ddot s+\omega_s^2s=0,\quad \omega_s^2=k/m;
\qquad \ddot d+\omega_a^2d=0,\quad \omega_a^2=(k+2\kappa)/m.
\]
两个方程的通解为$s=C_s\cos\omega_st+D_s\sin\omega_st$、$d=C_a\cos\omega_at+D_a\sin\omega_at$。初始$s(0)=d(0)=A/2$，给出$C_s=C_a=A/2$；求导后代入零初速度，得到$\omega_sD_s=\omega_aD_a=0$。因此$D_s=D_a=0$，还原$x_1=s+d$、$x_2=s-d$，得到
\[
\boxed{x_1=\frac A2(\cos\omega_st+\cos\omega_at),\quad
x_2=\frac A2(\cos\omega_st-\cos\omega_at).}
\]
用 $x_1=s+d$、$x_2=s-d$ 代回能量，动能中的交叉项相消，三根弹簧的势能化为
\[
T=m(\dot s^2+\dot d^2),\qquad U=ks^2+(k+2\kappa)d^2.
\]
因此s、d各自对应的等效质量为 $2m$。图\ref{sup4:fig:d7-e3-s}中的两条包络来自这两个独立简谐运动的叠加。
总能量写成 $E=m(\dot s^2+\omega_s^2s^2)+m(\dot d^2+\omega_a^2d^2)$，故
\[
E_s=\frac14kA^2,\qquad E_a=\frac14(k+2\kappa)A^2,
\quad E=\frac12(k+\kappa)A^2.
\]
两模态的初始位移系数均为 $A/2$，能量仍因频率不同而有差别。
\begin{figure}[H]\centering
\begin{tikzpicture}
\draw[axis](0,0)--(6.6,0)node[right]{$\tau=\Delta\omega t/2$};\draw[axis](0,-1.6)--(0,1.8)node[above]{包络系数};\draw[curve,domain=0:180,samples=90]plot({\x/30},{1.3*cos(\x)});\draw[guide,domain=0:180,samples=90]plot({\x/30},{1.3*sin(\x)});\node[below]at(3,0){$\pi/2$};\node[below]at(6,0){$\pi$};\node[left]at(0,1.3){$1$};\node[left]at(0,-1.3){$-1$};\node at(5.2,1.4){$x_2$：虚线};\node at(5,-1.8){$x_1$：实线};
\end{tikzpicture}
\caption{有符号慢包络；各自的振幅包络还须取绝对值。}\label{sup4:fig:d7-e3-s}
\end{figure}
令 $\bar\omega=(\omega_a+\omega_s)/2$、$\Delta\omega=\omega_a-\omega_s$，则
\[
x_1=A\cos\bar\omega t\cos\frac{\Delta\omega t}{2},\qquad
x_2=A\sin\bar\omega t\sin\frac{\Delta\omega t}{2}.
\]
图\ref{sup4:fig:d7-e3-s}给出两个慢因子。若$\kappa\ll k$，将$\omega_a=\omega_s\sqrt{1+2\kappa/k}$展开到一阶，得$\omega_a\simeq\omega_s(1+\kappa/k)$，所以$\Delta\omega\simeq\kappa/(m\omega_s)$。首次包络交接的时间为$t_b=\pi/\Delta\omega$。对$x_1$按乘积法则求导，
\[
\dot x_1=-A\bar\omega\sin\bar\omega t\cos\frac{\Delta\omega t}{2}
-\frac{A\Delta\omega}{2}\cos\bar\omega t\sin\frac{\Delta\omega t}{2}.
\]
代入$\Delta\omega t_b/2=\pi/2$，得到
\[
x_1(t_b)=0,\qquad
\dot x_1(t_b)=-\frac{A\Delta\omega}{2}\cos(\bar\omega t_b),
\]
第一块一般仍有非零速度。即使某次恰有 $x_1=\dot x_1=0$，中间弹簧也可能储能，见习题\ref{sup4:ex:d7-h4}。弱耦合时可近似说能量主要转移到另一物块，精确分配则须将三根弹簧全部计入。
\end{solution}



\begin{example}[失谐双摆与频率避交]\label{sup4:problem-example-5}\label{sup4:ex:d7-e5}
两单摆摆球质量均为 $m$，摆长为 $\ell_1,\ell_2$，支点高度使两球在静止时等高。两球由水平轻弹簧 $\kappa$ 连接，平衡时无形变。令球的水平小位移为 $x_1,x_2$。求小振动频率和振型；保持 $\omega_c^2=(g/\ell_1+g/\ell_2)/2$ 不变，改变失谐量 $d=(g/\ell_1-g/\ell_2)/2$，讨论两频率是否相交。
\end{example}
\begin{figure}[H]\centering
\begin{tikzpicture}
\fill(0,3)circle(1.6pt);\fill(4.5,2.3)circle(1.6pt);\draw[curve](0,3)--(0,0)node[midway,left]{$\ell_1$};\draw[curve](4.5,2.3)--(4.5,0)node[midway,right]{$\ell_2$};\fill(0,0)circle(3pt)node[below]{$m$};\fill(4.5,0)circle(3pt)node[below]{$m$};\draw[spring](.15,0)--(4.35,0)node[midway,above]{$\kappa$};\draw[axis](-.3,-.7)--(.8,-.7)node[right]{$x_1$};\draw[axis](4.2,-.7)--(5.3,-.7)node[right]{$x_2$};
\end{tikzpicture}
\caption{平衡时两摆球等高，支点高度不同。弹簧的一阶伸长为 $x_2-x_1$。}\label{sup4:fig:d7-e5-q}
\end{figure}
\begin{solution}
图\ref{sup4:fig:d7-e5-q}中，角度近似为 $x_i/\ell_i$。二次动能为 $m(\dot x_1^2+\dot x_2^2)/2$，势能为
\[
U_2=\frac{mg}{2\ell_1}x_1^2+\frac{mg}{2\ell_2}x_2^2
+\frac\kappa2(x_2-x_1)^2.
\]
势能的两个偏导分别为$\partial U_2/\partial x_1=mgx_1/\ell_1+\kappa(x_1-x_2)$、$\partial U_2/\partial x_2=mgx_2/\ell_2+\kappa(x_2-x_1)$。各令其与$m\ddot x_i$之和为零，记$c=\kappa/m>0$，得
\[
\ddot x_1+(\omega_c^2+d+c)x_1-cx_2=0,\qquad
\ddot x_2+(\omega_c^2-d+c)x_2-cx_1=0.
\]
代入$x_i=X_i\cos\omega t$，频率平方是矩阵
$\left(\begin{smallmatrix}\omega_c^2+d+c&-c\\-c&\omega_c^2-d+c\end{smallmatrix}\right)$ 的特征值，故
\[
\boxed{\omega_\pm^2=\omega_c^2+c\pm\sqrt{d^2+c^2}.}
\]
令 $\lambda=\omega^2$，特征行列式可先写成
\[
[(\omega_c^2+c)-\lambda]^2-d^2-c^2=0.
\]
这显示频率分支关于平均值 $\omega_c^2+c$ 对称。其两根之积为 $(g/\ell_1+c)(g/\ell_2+c)-c^2>0$，两根之和也为正，所以允许参数下两根均为正。
图\ref{sup4:fig:d7-e5-s}的纵轴是频率平方；实际频率差为 $2\sqrt{d^2+c^2}/(\omega_++\omega_-)$，不能把图上的间隔 $2c$ 当成频率差。
两摆长均为正还限制了 $|d|<\omega_c^2$。在允许范围内，两频率平方之差为 $2\sqrt{d^2+c^2}\ge2c$，非零耦合使两个分支保持分离。
\begin{figure}[H]\centering
\begin{tikzpicture}
\draw[axis](-3,0)--(3.2,0)node[right]{$d/c$};\draw[axis](0,0)--(0,4.7)node[above]{$\omega^2/c$};\draw[curve,domain=-2:2,samples=100]plot({1.3*\x},{.6*(5+sqrt((\x)^2+1))});\draw[curve,domain=-2:2,samples=100]plot({1.3*\x},{.6*(5-sqrt((\x)^2+1))});\draw[<->](.3,2.4)--(.3,3.6)node[midway,right]{$2$};\node[left]at(0,2.4){$4$};\node[left]at(0,3.6){$6$};\node at(-2,4.8){$\omega_c^2=4c$};
\end{tikzpicture}
\caption{有限耦合使频率平方分支避交；纵轴标出的间隔对应实际频率平方差 $2c$。}\label{sup4:fig:d7-e5-s}
\end{figure}
对低频模态，第一行方程给出
\[
\frac{X_2}{X_1}=\frac{d+\sqrt{d^2+c^2}}{c}>0,
\]
高频则为 $(d-\sqrt{d^2+c^2})/c<0$。在 $d=0$ 处恢复同相、反相振型。若 $d\gg c$ 且仍满足摆长允许范围，低频主要由第二摆运动；若 $d\ll-c$，低频主要由第一摆运动。图\ref{sup4:fig:d7-e5-s}的分支连续，跨过失谐零点后，低频模态中运动较强的摆球由第一摆换成第二摆。
\end{solution}



\begin{example}[双振子的反共振]\label{sup4:problem-example-8}\label{sup4:ex:d7-e8}
沿用例\ref{sup4:ex:d7-e3}的双振子，对第一块施加周期力 $F_0\cos\nu t$，第二块无外加驱动力。求无阻尼系统的周期特解，确定第一块发生反共振的频率及此时第二块的运动，讨论反共振与共振的差别。
\end{example}
\begin{figure}[H]\centering
\begin{tikzpicture}
\draw[curve](0,0)--(7.5,0);\draw[curve](0,0)--(0,1.2);\draw[curve](7.5,0)--(7.5,1.2);\draw[spring](0,.55)--(2,.55);\draw[curve](2,.2)rectangle(2.8,.9);\draw[spring](2.8,.55)--(4.6,.55);\draw[curve](4.6,.2)rectangle(5.4,.9);\draw[spring](5.4,.55)--(7.5,.55);\node at(2.4,.55){$m$};\node at(5,.55){$m$};\node[above]at(1,.7){$k$};\node[above]at(3.7,.7){$\kappa$};\node[above]at(6.4,.7){$k$};\draw[axis](2.4,1.3)--(4,1.3)node[right]{$F_0\cos\nu t$};
\end{tikzpicture}
\caption{驱动力只作用在第一块，两个位移正向均向右。}\label{sup4:fig:d7-e8-q}
\end{figure}
\begin{solution}
先在地面系列出受迫方程：
\[
\begin{aligned}
m\ddot x_1+(k+\kappa)x_1-\kappa x_2&=F_0\cos\nu t,\\
m\ddot x_2+(k+\kappa)x_2-\kappa x_1&=0.
\end{aligned}
\]
图\ref{sup4:fig:d7-e8-q}中外力只作用于第一块，却能通过中间弹簧驱动第二块。求同频周期特解时，两位移可取与外力同一余弦因子，但振幅系数允许为负，以表示反相。
图\ref{sup4:fig:d7-e8-q}中取 $x_i=X_i\cos\nu t$。令 $b=k+\kappa-m\nu^2$，幅值方程为
\[
\begin{pmatrix}b&-\kappa\\-\kappa&b\end{pmatrix}
\begin{pmatrix}X_1\\X_2\end{pmatrix}=
\begin{pmatrix}F_0\\0\end{pmatrix}.
\]
消元时不除以$b$，以免漏掉反共振点。将第一行乘$b$、第二行乘$\kappa$再相加，得到$(b^2-\kappa^2)X_1=bF_0$；将第一行乘$\kappa$、第二行乘$b$再相加，得到$(b^2-\kappa^2)X_2=\kappa F_0$。非共振时行列式$b^2-\kappa^2\ne0$，所以
\[
\boxed{X_1=\frac{F_0b}{b^2-\kappa^2},\qquad
X_2=\frac{F_0\kappa}{b^2-\kappa^2}.}
\]
分母在$\nu^2=k/m$、$(k+2\kappa)/m$为零，正是两固有频率。为求此时的响应，取$s=(x_1+x_2)/2$、$d=(x_1-x_2)/2$，相加、相减得到
\[
\ddot s+\omega_s^2s=\frac{F_0}{2m}\cos\nu t,\qquad
\ddot d+\omega_a^2d=\frac{F_0}{2m}\cos\nu t.
\]
若$\nu$等于其中一支固有频率，设该支的特解为$q_p=Ct\sin\nu t$。求导有
$\dot q_p=C\sin\nu t+C\nu t\cos\nu t$、
$\ddot q_p=2C\nu\cos\nu t-C\nu^2t\sin\nu t$，故
\[
\ddot q_p+\nu^2q_p=2C\nu\cos\nu t
=\frac{F_0}{2m}\cos\nu t,
\qquad q_p=\frac{F_0}{4m\nu}t\sin\nu t.
\]
它的振幅随时间增长，不能用有界周期特解代替；另一支模态仍按非共振方程求解。
\begin{figure}[H]\centering
\begin{tikzpicture}
\draw[axis](0,0)--(6.2,0)node[right]{$q=m\nu^2/k$};\draw[axis](0,-2.5)--(0,2.6)node[above]{$kX_1/F_0$};\draw[guide](1.4,-2.2)--(1.4,2.2);\draw[guide](4.2,-2.2)--(4.2,2.2);
\foreach\lo/\hi in{0/.85,1.15/2.85,3.15/4}{\draw[curve,domain=\lo:\hi,samples=100]plot({1.4*\x},{.55*(2-\x)/((1-\x)*(3-\x))});}
\fill(2.8,0)circle(1.7pt)node[below right,yshift=-5pt]{反共振};\node[below]at(1.4,0){$1$};\node[above left]at(2.8,0){$2$};\node[below]at(4.2,0){$3$};\node at(5.1,2.3){$\kappa=k$};\node[left]at(0,.55){$1$};
\end{tikzpicture}
\caption{第一块有符号周期响应。两个极点之间，响应在 $q=2$ 处过零。}\label{sup4:fig:d7-e8-s}
\end{figure}
图\ref{sup4:fig:d7-e8-s}的零点对应 $b=0$，此时
\[
\boxed{\nu_{\rm ar}^2=(k+\kappa)/m,\qquad X_1=0,\quad X_2=-F_0/\kappa.}
\]
此处 $b=0$ 而行列式为 $-\kappa^2\ne0$，因此反共振是响应分子的零点，不是固有频率处的奇点。由第二行方程 $-\kappa X_1=0$ 得 $X_1=0$，再由第一行 $-\kappa X_2=F_0$ 求出第二块振幅，既不需要也不允许令两块同时静止。
第一块保持在平衡位置，第二块的弹簧作用力 $\kappa x_2=-F_0\cos\nu t$ 恰好抵消驱动力。第二块则以第一块固定时的固有频率振动，故能维持这种周期运动。

这组周期解对应的初态为 $x_1=0,x_2=-F_0/\kappa$、两速度为零。从两块均静止于平衡位置开始驱动，还须叠加自由振动，详见习题\ref{sup4:ex:d7-h9}。无阻尼、非共振周期解的平均输入功率为零；加入阻尼后，响应幅值与耗散功率随之改变。
\end{solution}



\section{零频模态与线性近似的边界}

二次势能为零的方向，可能对应整体平移，也可能仍有高阶回复力。零频率的来源决定线性化是否足够，以及是否需要保留非线性项。



\subsection{零频率与小振动近似的范围}
$K$ 正定时，所有 $\omega_r^2>0$，线性化后的小振动有界。若 $\omega_r^2=-\sigma^2<0$，特征方程为 $\lambda^2-\sigma^2=0$，通解为 $Q_r=A\ee^{\sigma t}+B\ee^{-\sigma t}$；存在增长分支，原平衡不稳定。零特征值须直接解 $\ddot Q_0=0$：先积分得 $\dot Q_0=\dot Q_0(0)$，再积分得 $Q_0=Q_0(0)+\dot Q_0(0)t$。非零初速度产生漂移，不能将零频率直接代入含 $1/\omega$ 的简谐运动公式。

展开动能和势能时，须采用同一小量阶数。例如 $\rho\theta^2$、$\rho\dot\theta^2$ 通常为三阶能量项，在线性方程中可舍去，在更高阶运动中却会产生耦合。振动解的适用范围还受接触、纯滚动及弹簧线性等条件限制，各题中另行讨论。


\subsection{零频模态与四次回复势}
当势能沿某方向保持不变，零频模态给出匀速漂移。若二次项恰好消失，而最低非零项为$U-U_* =\beta q^4/4$，$\beta>0$，则运动方程变为
\[
 M_*\ddot q+\beta q^3=0.
\]
从振幅 $A>0$ 处静止释放，能量为 $\beta A^4/4$，因此
\[
\frac12M_*\dot q^2=\frac\beta4(A^4-q^4),\qquad
|\dot q|=\sqrt{\frac\beta{2M_*}}\sqrt{A^4-q^4}.
\]
利用四分之一周期的对称性，并令 $q=As$、$\dd q=A\dd s$，上下限 $q=0,A$ 对应 $s=0,1$，得到
\[
\frac T4=\sqrt{\frac{2M_*}\beta}\int_0^A\frac{\dd q}{\sqrt{A^4-q^4}}
=\frac1A\sqrt{\frac{2M_*}\beta}\int_0^1\frac{\dd s}{\sqrt{1-s^4}}.
\]
最后的积分是与振幅无关的有限常数，不能用初等函数给出原函数。于是
\[
 T(A)=\frac4A\sqrt{\frac{2M_*}{\beta}}
 \int_0^1\frac{\dd s}{\sqrt{1-s^4}}.
\]
回复运动的周期随振幅而变；$A\to0$时周期趋于无穷，故不能赋予它非零的线性固有频率。例\ref{ex:hw4-4}的临界木条和例\ref{ex:hw4-8}的两个端点电荷比，都须保留四次势能才能判断。若多坐标系统还有非退化方向，应同时考虑耦合高阶项，不能只沿一条坐标轴画势能曲线就结束稳定性分析。



\LectureNeedspace{4\baselineskip}


\subsection{例题与模型变式}

\begin{example}[自由双质点的零频模态]\label{sup4:problem-example-10}\label{sup4:ex:d7-e10}
光滑直线上，质量为 $m_1,m_2$ 的两个物体由劲度系数为 $\kappa$、原长为 $\ell_0$ 的弹簧连接，两端均不固定。初始弹簧无形变、两物体静止，对第一物体施加向右冲量 $J$。求两个物体相对初始位置的位移及能量分配，说明零频模态的含义。
\end{example}
\begin{figure}[H]\centering
\begin{tikzpicture}
\draw[curve](-.5,0)--(7.5,0);\draw[curve](0,.2)rectangle(1,.9);\draw[curve](5.5,.2)rectangle(6.9,.9);\draw[spring](1,.55)--(5.5,.55)node[midway,above]{$\kappa,\ \ell_0$};\node at(.5,.55){$m_1$};\node at(6.2,.55){$m_2$};\draw[axis](.5,1.3)--(2,1.3)node[right]{$J$};\draw[axis](2,-.6)--(3.3,-.6)node[right]{位移正向};
\end{tikzpicture}
\caption{两端没有固定墙，弹簧只改变相对运动，不限制整体平移。}\label{sup4:fig:d7-e10-q}
\end{figure}
\begin{solution}
令总质量 $M=m_1+m_2$、约化质量 $\mu=m_1m_2/M$。在图\ref{sup4:fig:d7-e10-q}中，以质心位置 $X$ 和弹簧伸长 $r=x_2-x_1-\ell_0$ 为坐标，有
\[
T=\frac12M\dot X^2+\frac12\mu\dot r^2,\qquad U=\frac12\kappa r^2.
\]
$x_1,x_2$ 表示地面位置，所求位移为各自相对初始位置的 $u_1,u_2$。反解坐标变换，
\[
x_1=X-\frac{m_2}{M}(\ell_0+r),\qquad
x_2=X+\frac{m_1}{M}(\ell_0+r).
\]
对时间求导再代入两质点动能，交叉项系数为 $-m_1m_2/M+m_2m_1/M=0$，径向项的系数为约化质量 $\mu$。这解释了质心平移与弹簧伸缩为何能够分离。
方程完全分离为 $\ddot X=0$、$\ddot r+\omega^2r=0$，其中
\[
\omega^2=\frac\kappa\mu=\kappa\left(\frac1{m_1}+\frac1{m_2}\right).
\]
质心方程积分两次，通解为$X=C_0+C_1t$；相对方程的特征根为$\pm\mathrm i\omega$，通解为$r=A\cos\omega t+B\sin\omega t$。冲量后$X(0)=X_0$、$\dot X(0)=J/M$，所以$C_0=X_0,C_1=J/M$；又$r(0)=0$给出$A=0$，$\dot r(0)=\omega B=-J/m_1$给出$B=-J/(m_1\omega)$。因此
\[
X-X_0=\frac JM t,\qquad r=-\frac{J}{m_1\omega}\sin\omega t.
\]
\begin{figure}[H]\centering
\begin{tikzpicture}
\begin{scope}
\draw[axis](0,0)--(3.8,0)node[right]{$t$};\draw[axis](0,0)--(0,2.2)node[above]{$X-X_0$};\draw[curve](0,0)--(3.2,1.7)node[right]{$Jt/M$};\node at(1.5,-.6){整体平移};
\end{scope}
\begin{scope}[xshift=5.4cm]
\draw[axis](0,0)--(3.8,0)node[right]{$\omega t$};\draw[axis](0,-1.4)--(0,1.4)node[above]{$r$};\draw[curve,domain=0:360,samples=90]plot({\x/110},{-sin(\x)});\node at(1.6,-1.9){内部振动};
\end{scope}
\end{tikzpicture}
\caption{质心匀速漂移，弹簧伸长保持有界。右图纵向以 $J/(m_1\omega)$ 为单位。}\label{sup4:fig:d7-e10-s}
\end{figure}
由质心分配关系，两物体相对初始位置的位移为
\[
\boxed{u_1=\frac JM t+\frac{m_2J}{Mm_1\omega}\sin\omega t,
\qquad u_2=\frac JM t-\frac{J}{M\omega}\sin\omega t.}
\]
图\ref{sup4:fig:d7-e10-s}中的零频分支对应同向平移振型 $(1,1)$，其时间因子为常量加线性项。扣除这一漂移后，内部振动的两位移满足质量加权和为零。

能量分为
\[
E_{\rm cm}=\frac{J^2}{2M},\qquad
E_{\rm int}=\frac12\mu\frac{J^2}{m_1^2}
=\frac{m_2J^2}{2m_1M},
\]
两部分相加为 $J^2/(2m_1)$。内部运动保持有界，两物体的地面位置则随质心漂移；讨论内部平衡的稳定性时，应先扣除整体平移。
\end{solution}



\begin{example}[弹簧摆的线性分离与内共振]\label{deep4:problem-example-7}\label{deep4:ex:d7-e7}
一轻弹簧劲度系数为 $k$、原长为 $\ell_0>0$，上端固定，下端连接质量为 $m$ 的小球，小球可在竖直平面运动。取弹簧长度 $r$、偏离竖直向下的角度 $\theta$ 为坐标。求平衡位置附近的径向、摆动频率，并求两者满足 $\omega_r=2\omega_\theta$ 时的参数条件。再令首阶摆动为 $\theta=\Theta\cos\omega_\theta t$，求它引起的最低阶径向驱动项。
\end{example}
\begin{figure}[H]\centering
\begin{tikzpicture}
\fill(0,0)circle(1.6pt)node[above]{$O$};\draw[guide](0,0)--(0,-3.6);\draw[spring](0,0)--(1,-3);\fill(1,-3)circle(3pt)node[right]{$m$};\node[right]at(.7,-1.6){$r$};\draw(0,-.85)arc(-90:-72:.85);\node at(.18,-1.6){$\theta$};\node[align=left]at(4,-1.4){$k,\ \ell_0$\\平衡长度为 $r_0$};
\end{tikzpicture}
\caption{径向伸缩与角向摆动使用不同坐标；平衡态为竖直悬垂。}\label{deep4:fig:d7-e7-q}
\end{figure}
\begin{solution}
图\ref{deep4:fig:d7-e7-q}中的精确能量为
\[
T=\frac12m(\dot r^2+r^2\dot\theta^2),\qquad
U=\frac12k(r-\ell_0)^2-mgr\cos\theta.
\]
先列出能量中的各项导数：
\[
\frac{\partial U}{\partial r}=k(r-\ell_0)-mg\cos\theta,
\quad\frac{\partial U}{\partial\theta}=mgr\sin\theta,
\quad\frac{\partial T}{\partial r}=mr\dot\theta^2,
\]
\[
\frac{\dd}{\dd t}\frac{\partial T}{\partial\dot r}=m\ddot r,
\qquad
\frac{\dd}{\dd t}\frac{\partial T}{\partial\dot\theta}
=\frac{\dd}{\dd t}(mr^2\dot\theta)
=mr^2\ddot\theta+2mr\dot r\dot\theta.
\]
平衡时令两个势能偏导为零，取向下悬垂分支，得$r_0=\ell_0+mg/k$、$\theta=0$。置 $\rho=r-r_0$，按 $\rho/r_0$ 与 $\theta$ 同为小量展开，二次项为
\[
T_2=\frac12m\dot\rho^2+\frac12mr_0^2\dot\theta^2,
\quad U_2=U_*+\frac12k\rho^2+\frac12mgr_0\theta^2.
\]
例如动能中的 $mr_0\rho\dot\theta^2$、势能中的 $mg\rho\theta^2/2$ 都是三阶量，在线性方程所需的二次能量中应舍去。径向一次势能项则由 $k(r_0-\ell_0)=mg$ 抵消。两个方向在线性阶分离，并不表示精确运动没有耦合。
对二次能量求导，径向方程为$m\ddot\rho+k\rho=0$，角向方程为$mr_0^2\ddot\theta+mgr_0\theta=0$。分别除以各自的惯性系数，两个线性方程独立，其频率为
\[
\boxed{\omega_r=\sqrt{k/m},\qquad \omega_\theta=\sqrt{g/r_0}.}
\]
条件 $\omega_r^2=4\omega_\theta^2$ 等价于 $kr_0=4mg$，利用平衡长度，得到
\[
\boxed{k\ell_0=3mg.}
\]
\begin{figure}[H]\centering
\begin{tikzpicture}
\begin{scope}
\draw[axis](-.5,0)--(4,0)node[right]{驱动频率};\draw[axis](-.5,-1.5)--(-.5,1.5)node[above]{二阶径向驱动};\draw[curve](0,0)--(0,.4);\node[below]at(0,0){$0$};\draw[curve](2.5,0)--(2.5,-1.2);\node[above]at(2.5,0){$2\omega_\theta$};\node[right]at(0,.45){$g\Theta^2/4$};\node[right]at(2.5,-1.2){$-3g\Theta^2/4$};
\end{scope}
\end{tikzpicture}
\caption{角向运动产生恒定项和倍频项；竖线表示傅里叶系数及符号。}\label{deep4:fig:d7-e7-s}
\end{figure}
精确径向方程为 $\ddot r-r\dot\theta^2=g\cos\theta-(k/m)(r-\ell_0)$。保留 $\Theta^2$ 阶，得
\[
\ddot\rho+\omega_r^2\rho
=r_0\dot\theta^2-\frac g2\theta^2
=\frac{g\Theta^2}{4}-\frac{3g\Theta^2}{4}\cos2\omega_\theta t.
\]
代入时，$r_0\omega_\theta^2=g$，而 $\sin^2\omega_\theta t=(1-\cos2\omega_\theta t)/2$、$\cos^2\omega_\theta t=(1+\cos2\omega_\theta t)/2$。离心项与重力二次项共同给出常量 $g\Theta^2/4$ 和倍频项 $-3g\Theta^2\cos2\omega_\theta t/4$，两项都应保留。常量只平移径向振动中心，倍频项才可能发生共振。
图\ref{deep4:fig:d7-e7-s}中，第二根谱线在$2\omega_\theta=\omega_r$时与径向固有频率相合。常量驱动的特解为$g\Theta^2/(4\omega_r^2)$。对余下的倍频驱动，设$\rho_{\rm sec}=Ct\sin\omega_rt$，逐次求导：
\[
\dot\rho_{\rm sec}=C\sin\omega_rt+C\omega_rt\cos\omega_rt,
\qquad
\ddot\rho_{\rm sec}=2C\omega_r\cos\omega_rt-C\omega_r^2t\sin\omega_rt.
\]
代入左端后含$t\sin\omega_rt$的两项抵消，故$2C\omega_r=-3g\Theta^2/4$，得到
\[
\rho_{\rm sec}=-\frac{3g\Theta^2}{8\omega_r}t\sin\omega_rt.
\]
在这一阶近似内，径向通解还包括$A\cos\omega_rt+B\sin\omega_rt$及上述常量特解，$A,B$由径向初态确定。用前述角坐标的导数代入拉格朗日方程，并除以$mr$，精确角向方程为
\[
r\ddot\theta+2\dot r\dot\theta+g\sin\theta=0.
\]
径向变化会通过这两个含r的项反过来影响摆角；图\ref{deep4:fig:d7-e7-s}中的“驱动”是同一系统内部的能量交换。
随时间增长的项说明首阶角向解只适用于早期运动。实际系统总能量守恒，径向能量增加会反过来改变摆动；当 $|\rho|/r_0$ 或角向振幅的变化已不可忽略时，须将两个非线性方程联立求解。

还可以由能量估计早期近似的时间范围。在精确内共振$\omega_r=2\omega_\theta$处，利用$g=r_0\omega_\theta^2$，久期项的振幅为
\[
 \frac{A_\rho(t)}{r_0}=\frac{3\Theta^2\omega_\theta t}{16}.
\]
取多个振动周期以后的缓变包络，径向振动能约为$m\omega_r^2A_\rho^2/2$，角向初始能量约为$mr_0^2\omega_\theta^2\Theta^2/2$，两者之比为
\[
 \frac{E_r}{E_{\theta,0}}\simeq
 \frac{4A_\rho^2}{r_0^2\Theta^2}
 =\frac9{64}\Theta^2(\omega_\theta t)^2.
\]
忽略角向振幅的变化至少要求$\Theta\omega_\theta t\ll1$。当$t$达到$1/(\Theta\omega_\theta)$量级，能量交换已不能视为微小修正，即使径向位移相对于$r_0$仍小，固定角向振幅的驱动近似也会失效。此估计给出近似失效的尺度，并不预测能量随后继续无界增长。
\end{solution}



\section{习题与解答}

\begin{exercise}[双摆下端受冲量]\label{sup4:problem-exercise-7}\label{sup4:ex:d7-h7}
例\ref{sup4:ex:d7-e6}的等长等质量双摆初始竖直静止，对下端质点施加向右水平冲量 $J$。求冲量后的两个角速度、此后的小振动及各模态能量。设冲量足够小，使随后摆角均可线性化。
\end{exercise}
\begin{figure}[H]\centering
\begin{tikzpicture}
\fill(0,0)circle(1.6pt)node[above]{$O$};\draw[curve](0,0)--(0,-3.6);\fill(0,-1.8)circle(2.5pt)node[left]{$m$};\fill(0,-3.6)circle(2.5pt)node[left]{$m$};\node[right]at(0,-.9){$\ell$};\node[right]at(0,-2.7){$\ell$};\draw[axis](0,-3.6)--(1.8,-3.6)node[right]{$J$};\node[align=left]at(3,-1.4){两角从竖直向下量起\\正向均向右};
\end{tikzpicture}
\caption{冲量作用在下端质点，两个角坐标的广义冲量都须由作用点位移求出。}\label{sup4:fig:d7-h7-q}
\end{figure}
\begin{solution}
图\ref{sup4:fig:d7-h7-q}中，下端水平位置为 $x_B=\ell(\sin\theta_1+\sin\theta_2)$。在竖直构形，虚位移为 $\delta x_B=\ell(\delta\theta_1+\delta\theta_2)$，故广义冲量为 $(J\ell,J\ell)$。利用双摆的质量矩阵，
\[
m\ell^2\begin{pmatrix}2&1\\1&1\end{pmatrix}
\begin{pmatrix}\dot\theta_1(0^+)\\\dot\theta_2(0^+)\end{pmatrix}
=J\ell\begin{pmatrix}1\\1\end{pmatrix},
\]
竖直构形处，质量矩阵的无量纲部分逆矩阵为
\[
\begin{pmatrix}2&1\\1&1\end{pmatrix}^{-1}
=\begin{pmatrix}1&-1\\-1&2\end{pmatrix}.
\]
乘上广义冲量后，第一行两个分量抵消，第二行剩一个单位分量。这一消元说明实际施加在下端的冲量，不能只记在第二个角坐标上。
解得 $\dot\theta_1(0^+)=0$、$\dot\theta_2(0^+)=J/(m\ell)$。上端质点初始水平速度仍为零，下端速度为 $J/m$，与初始两杆均竖直、约束冲量不产生水平分量相符。
\begin{figure}[H]\centering
\begin{tikzpicture}
\node at(1.2,2.3){冲量后瞬间};\fill(0,1)circle(2pt)node[left]{上端};\fill(0,-.4)circle(2pt)node[left]{下端};\node[right]at(0,1){$v_A=0$};\draw[axis](0,-.4)--(1.6,-.4)node[right]{$v_B=J/m$};
\draw[axis](4,0)--(7,0);\draw[axis](4,0)--(4,2.4)node[above]{能量比例};\draw[curve,fill=gray!15](4.4,0)rectangle(5.3,1.707);\draw[curve,fill=gray!15](6,0)rectangle(6.9,.293);\node[above]at(4.85,1.707){$(2+\sqrt2)/4$};\node[above]at(6.45,.35){$(2-\sqrt2)/4$};\node[below]at(4.85,0){低频};\node[below]at(6.45,0){高频};
\end{tikzpicture}
\caption{初始水平速度及两模态的能量份额，低频模态获得较多能量。}\label{sup4:fig:d7-h7-s}
\end{figure}
两振型仍为 $(1,\sqrt2)$、$(1,-\sqrt2)$。由初始位移为零，两模态时间因子均为正弦。先求模态初速度 $b_-,b_+$：由图\ref{sup4:fig:d7-h7-s}的实际初速度，
\[
b_-+b_+=0,\qquad
\sqrt2(b_--b_+)=\frac{J}{m\ell}.
\]
第一式给 $b_+=-b_-$，代入第二式得 $2\sqrt2b_-=J/(m\ell)$，所以
\[
b_-=\frac{J}{2\sqrt2m\ell},\qquad b_+=-\frac{J}{2\sqrt2m\ell}.
\]
模态坐标分别为 $q_\pm=(b_\pm/\omega_\pm)\sin\omega_\pm t$；再由 $\theta_1=q_-+q_+$、$\theta_2=\sqrt2(q_--q_+)$，得到
\[
\boxed{\begin{aligned}
\theta_1&=\frac{J}{2\sqrt2m\ell}\left(\frac{\sin\omega_-t}{\omega_-}-\frac{\sin\omega_+t}{\omega_+}\right),\\
\theta_2&=\frac{J}{2m\ell}\left(\frac{\sin\omega_-t}{\omega_-}+\frac{\sin\omega_+t}{\omega_+}\right).
\end{aligned}}
\]
未归一化振型的等效惯量为 $m\ell^2(4+2\sqrt2)$、$m\ell^2(4-2\sqrt2)$，两模态初始速度系数的大小均为 $J/(2\sqrt2m\ell)$。因此
\[
E_- =\frac{J^2}{8m}(2+\sqrt2),\qquad
E_+ =\frac{J^2}{8m}(2-\sqrt2).
\]
相加为 $J^2/(2m)$，等于图\ref{sup4:fig:d7-h7-s}中下端质点的初始动能。若只给第二角坐标一个广义冲量，会错误地算出上端质点也获得水平初速度。
\end{solution}


\begin{exercise}[弹簧能否稳定两个倒立摆（选做）]\label{sup4:problem-exercise-12}\label{sup4:ex:d7-h12}
两个相同倒立摆各由轻杆和端点质量 $m$ 组成，杆长为 $\ell$，下端支点固定且等高，间距为 $D$。两球在竖直向上的平衡态由水平弹簧 $\kappa$ 连接，弹簧原长为 $D$，可伸可压。两杆可独立在同一竖直平面转动，摆角均从竖直向上向右量起。求线性模态，判断增大 $\kappa$ 能否使整个平衡稳定，再讨论严格反对称运动在临界刚度处的稳定性。
\end{exercise}
\begin{figure}[H]\centering
\begin{tikzpicture}
\fill(0,0)circle(1.7pt)node[below]{$P_1$};\fill(4.5,0)circle(1.7pt)node[below]{$P_2$};\draw[curve](0,0)--(0,2.5);\draw[curve](4.5,0)--(4.5,2.5);\fill(0,2.5)circle(2.7pt)node[above]{$m$};\fill(4.5,2.5)circle(2.7pt)node[above]{$m$};\draw[spring](.15,2.5)--(4.35,2.5)node[midway,above]{$\kappa$};\node[left]at(0,1.2){$\ell$};\node[right]at(4.5,1.2){$\ell$};\draw[<->](0,-.65)--(4.5,-.65)node[midway,below]{$D$};
\end{tikzpicture}
\caption{竖直向上的构形是相对位移零点，两杆均允许独立向左右偏转。}\label{sup4:fig:d7-h12-q}
\end{figure}
\begin{solution}
图\ref{sup4:fig:d7-h12-q}中，球的水平位移一阶为 $\ell\theta_i$，弹簧伸长为 $\ell(\theta_2-\theta_1)$。因此
\[
T_2=\frac12m\ell^2(\dot\theta_1^2+\dot\theta_2^2),\quad
U_2=U_*-\frac12mg\ell(\theta_1^2+\theta_2^2)
+\frac12\kappa\ell^2(\theta_2-\theta_1)^2.
\]
取同相、反相坐标 $u=(\theta_1+\theta_2)/2$、$v=(\theta_1-\theta_2)/2$，得到
\[
\ddot u-\frac g\ell u=0,\qquad
\ddot v+\left(\frac{2\kappa}{m}-\frac g\ell\right)v=0.
\]
\begin{figure}[H]\centering
\begin{tikzpicture}
\begin{scope}
\foreach\x in{0,2}{\fill(\x,0)circle(1.5pt);\draw[curve](\x,0)--++(.5,1.8);\fill(\x+.5,1.8)circle(2pt);}\draw[spring](.5,1.8)--(2.5,1.8);\draw[axis](.5,2.2)--(1.3,2.2);\draw[axis](2.5,2.2)--(3.3,2.2);\node at(1.3,-.65){同相：弹簧不变长};
\end{scope}
\begin{scope}[xshift=5cm]
\fill(0,0)circle(1.5pt);\fill(2.5,0)circle(1.5pt);\draw[curve](0,0)--(.45,1.8);\draw[curve](2.5,0)--(2.05,1.8);\fill(.45,1.8)circle(2pt);\fill(2.05,1.8)circle(2pt);\draw[spring](.45,1.8)--(2.05,1.8);\node at(1.25,-.65){反相：弹簧提供回复项};
\end{scope}
\end{tikzpicture}
\caption{同相倾斜保留原来的球间距离，耦合弹簧无法抑制该不稳定方向。}\label{sup4:fig:d7-h12-s}
\end{figure}
图\ref{sup4:fig:d7-h12-s}中的同相模态始终有$\omega_u^2=-g/\ell<0$。令$\gamma=\sqrt{g/\ell}$，特征方程$s^2-\gamma^2=0$有两个实根$\pm\gamma$，所以$u=C_+\ee^{\gamma t}+C_-\ee^{-\gamma t}$。例如初始$u(0)=\varepsilon,\dot u(0)=0$，则$C_++C_-=\varepsilon$、$\gamma(C_+-C_-)=0$，解得$C_+=C_-=\varepsilon/2$，即$u=\varepsilon\cosh(\gamma t)$。任意小偏差均会增长，所以任何有限$\kappa$都不能使整个平衡稳定。

同相运动时，$\theta_1=\theta_2$，两球具有相同的水平、竖直位移，相对位置始终为 $(D,0)$，弹簧长度保持D。图\ref{sup4:fig:d7-h12-s}左图的共同偏转因此始终不受弹簧的回复作用，这一结论在有限小偏角下也成立。
严格限定 $u=0$ 的反对称运动则在 $\kappa>mg/(2\ell)$ 时具有正线性频率平方。在临界 $\kappa=mg/(2\ell)$，线性频率为零，须继续展开势能。令 $\theta_1=v,\theta_2=-v$，小角范围内两球等高，弹簧伸长为 $-2\ell\sin v$，故
\[
U(v)=2mg\ell\cos v+2\kappa\ell^2\sin^2v.
\]
在临界刚度处，第二项系数为$mg\ell$。将$\cos v=1-v^2/2+v^4/24+O(v^6)$、$\sin^2v=v^2-v^4/3+O(v^6)$分别代入，得
\[
\begin{aligned}
U(v)&=2mg\ell-mg\ell v^2+\frac{mg\ell}{12}v^4
+mg\ell v^2-\frac{mg\ell}{3}v^4+O(v^6)\\
&=U_*-\frac{mg\ell}{4}v^4+O(v^6).
\end{aligned}
\]
四次项为负，反对称子空间内的平衡也不稳定。即使刚度超过此临界值、反相模态转为稳定，同相模态的负频率平方仍使整个平衡不稳定。
在临界刚度、严格反对称的子空间内，动能为 $m\ell^2\dot v^2$，等效惯量为 $2m\ell^2$。由四次势能项求导，
\[
2m\ell^2\ddot v=-U'(v)=mg\ell v^3+O(v^5),
\qquad \ddot v=\frac g{2\ell}v^3+O(v^5).
\]
小的正偏角获得正加速度，小的负偏角获得负加速度，运动都远离平衡。它与第六讲临界圆环的正四次势阱符号相反，零线性频率本身不足以给出稳定性结论。
\end{solution}


\begin{exercise}[包络交接时能量在哪里]\label{sup4:problem-exercise-4}\label{sup4:ex:d7-h4}
例\ref{sup4:ex:d7-e3}中取 $\kappa=3k/2$，仍从 $x_1=A,x_2=0$、速度均为零释放。求 $t_* =\pi/\omega_s$ 时两块的位置、速度及三根弹簧各自储能。是否可以说此时全部能量都在第二块及其外侧弹簧中？
\end{exercise}
\begin{figure}[H]\centering
\begin{tikzpicture}
\draw[curve](0,0)--(7,0);\foreach\x in{0,7}{\draw[curve](\x,0)--(\x,1.1);}\draw[spring](0,.5)--(2,.5);\draw[curve](2,.2)rectangle(2.7,.8);\draw[spring](2.7,.5)--(4.2,.5);\draw[curve](4.2,.2)rectangle(4.9,.8);\draw[spring](4.9,.5)--(7,.5);\node[above]at(1,.6){$k$};\node[above]at(3.45,.6){$3k/2$};\node[above]at(5.9,.6){$k$};\node[below]at(2.35,0){初始 $A$};\node[below]at(4.55,0){初始 $0$};
\end{tikzpicture}
\caption{耦合不弱，中间弹簧的储能不能忽略。}\label{sup4:fig:d7-h4-q}
\end{figure}
\begin{solution}
图\ref{sup4:fig:d7-h4-q}的反相频率为 $\omega_a=\sqrt{(k+3k)/m}=2\omega_s$。在 $t_* =\pi/\omega_s$，有 $\cos\omega_st_*=-1$、$\cos\omega_at_*=1$，两正弦均为零，所以
\[
x_1(t_*)=0,\qquad x_2(t_*)=-A,\qquad
\dot x_1(t_*)=\dot x_2(t_*)=0.
\]
此时左弹簧无形变，右弹簧形变量大小为 $A$，中间弹簧形变量为 $x_2-x_1=-A$。
\begin{figure}[H]\centering
\begin{tikzpicture}
\draw[axis](0,0)--(6.5,0)node[right]{储能部分};\draw[axis](0,0)--(0,2.7)node[above]{$E_i/E$};\draw[curve](1,0)rectangle(2,0);\draw[curve,fill=gray!15](3,0)rectangle(4,2.1);\draw[curve,fill=gray!15](5,0)rectangle(6,1.4);\node[below]at(1.5,0){左弹簧};\node[below]at(3.5,0){中间弹簧};\node[below]at(5.5,0){右弹簧};\node[above]at(1.5,0){$0$};\node[above]at(3.5,2.1){$3/5$};\node[above]at(5.5,1.4){$2/5$};
\end{tikzpicture}
\caption{指定时刻的能量分布。两物块动能均为零，中间弹簧仍储存五分之三的总能量。}\label{sup4:fig:d7-h4-s}
\end{figure}
由图\ref{sup4:fig:d7-h4-s}，
\[
U_L=0,\qquad U_C=\frac12\kappa A^2=\frac34kA^2,\qquad U_R=\frac12kA^2.
\]
总能量为 $5kA^2/4$，等于初态的 $\frac12kA^2+\frac12\kappa A^2$。第一块虽瞬时静止于平衡位置，右外侧弹簧中仍只有总能量的五分之二，其余五分之三储存在耦合弹簧中，不能全归于某一块独立振子的能量。
瞬时静止也不表示此后停在原处。由第一块受力方程，在 $t_*$ 有
\[
\ddot x_1(t_*)=\frac{\kappa x_2(t_*)}{m}=-\frac{\kappa A}{m}\ne0.
\]
耦合弹簧的力随后使两块重新运动。图\ref{sup4:fig:d7-h4-s}记录的是一个瞬时能量分布；两速度及弹簧能量变化率在此刻为零，受力却未平衡。
\end{solution}


\begin{exercise}[反共振频率下从静止启动]\label{sup4:problem-exercise-9}\label{sup4:ex:d7-h9}
例\ref{sup4:ex:d7-e8}中取 $\kappa=k$，从两块均静止于平衡位置开始，对第一块施加 $F_0\cos\nu t$，其中 $\nu=\sqrt{2k/m}$。求完整运动，并解释为什么第一块仍会运动。
\end{exercise}
\begin{figure}[H]\centering
\begin{tikzpicture}
\draw[axis](0,0)--(6.2,0)node[right]{频率};\draw[axis](0,0)--(0,1.9);\draw[curve](1.3,0)--(1.3,.8);\draw[curve](3,0)--(3,1.4);\draw[curve](4.8,0)--(4.8,.8);\node[below]at(1.3,0){$\omega_s$};\node[below]at(3,0){$\nu$};\node[below]at(4.8,0){$\omega_a$};\node[above]at(3,1.4){驱动频率};\node at(3,-1){初态：$x_1=x_2=\dot x_1=\dot x_2=0$};
\end{tikzpicture}
\caption{驱动频率位于两固有频率之间，频率位置仅按次序示意。}\label{sup4:fig:d7-h9-q}
\end{figure}
\begin{solution}
图\ref{sup4:fig:d7-h9-q}中 $\omega_s=\sqrt{k/m}$、$\omega_a=\sqrt{3k/m}$。反共振周期特解为 $x_{1p}=0$、$x_{2p}=-(F_0/k)\cos\nu t$。它在 $t=0$ 的第二块位移不为零，因此尚未满足题设初态。

令$s=(x_1+x_2)/2$、$d=(x_1-x_2)/2$，两个模态方程均具有形式
\[
\ddot q+\omega^2q=\frac{F_0}{2m}\cos\nu t,
\qquad q(0)=\dot q(0)=0,
\]
其中$\omega$分别取$\omega_s,\omega_a$。由于$\nu\ne\omega$，设特解$q_p=C\cos\nu t$，代入得$C(\omega^2-\nu^2)=F_0/(2m)$。通解是
\[
q=A\cos\omega t+B\sin\omega t+C\cos\nu t.
\]
在$t=0$代入位移条件，得$A+C=0$；求导并代入速度条件，得$B\omega=0$。所以$A=-C,B=0$，两模态的完整解为
\[
\begin{aligned}
s&=\frac{F_0}{2m(\omega_s^2-\nu^2)}(\cos\nu t-\cos\omega_st),\\
d&=\frac{F_0}{2m(\omega_a^2-\nu^2)}(\cos\nu t-\cos\omega_at).
\end{aligned}
\]
本题两个分母分别为 $-2k$、$2k$，相加求 $x_1=s+d$ 时，驱动频率的余弦项恰好相消；固有频率项则继续保留。这与图\ref{sup4:fig:d7-h9-s}的三条频率标记对应。
自由解在初始时须补入位移 $(0,F_0/k)$，初速度仍为零。将这一位移分解成同相与反相部分，得到
\[
\boxed{\begin{aligned}
x_1&=\frac{F_0}{2k}(\cos\omega_st-\cos\omega_at),\\
x_2&=-\frac{F_0}{k}\cos\nu t+\frac{F_0}{2k}(\cos\omega_st+\cos\omega_at).
\end{aligned}}
\]
\begin{figure}[H]\centering
\begin{tikzpicture}
\draw[axis](0,0)--(6.2,0)node[right]{频率};\draw[axis](0,-1.5)--(0,1.7)node[above]{第一块余弦系数};\draw[curve](1.3,0)--(1.3,1);\fill(3,0)circle(1.7pt);\draw[curve](4.8,0)--(4.8,-1);\node[above]at(1.3,1){$F_0/(2k)$};\node[below]at(4.8,-1){$-F_0/(2k)$};\node[below]at(1.3,0){$\omega_s$};\node[above]at(3,0){$0$};\node[below]at(3,-.35){$\nu$};\node[above]at(4.8,0){$\omega_a$};
\end{tikzpicture}
\caption{第一块在驱动频率处的系数为零，但两个自由频率分量均不为零。}\label{sup4:fig:d7-h9-s}
\end{figure}
图\ref{sup4:fig:d7-h9-s}中，第一块的 $\nu$ 分量为零，两个自由频率分量则因初始条件而保留。无阻尼时它们不衰减，这就是启动后第一块继续运动的原因。代入 $t=0$，得 $x_1=0$、$\ddot x_1=F_0/m$，与初始外力相符。
\end{solution}


\begin{exercise}[由频率求弹簧与模态能量]\label{deep4:problem-exercise-3}\label{deep4:ex:d7-h3}
例\ref{sup4:ex:d7-e3}的对称双振子已测得同相、反相频率分别为 $\omega_s,\omega_a$，且 $\omega_a>\omega_s$，两质量 $m$ 已知。求 $k,\kappa$。若初始 $x_1=A>0,x_2=B$，$0<B<A$，两速度为零，求使两个模态能量相等的 $B$。
\end{exercise}
\begin{figure}[H]\centering
\begin{tikzpicture}
\foreach\y/\txt in{1/同相,-1/反相}{\node[left]at(0,\y){\txt};\fill(1,\y)circle(2.3pt);\fill(4,\y)circle(2.3pt);}\draw[axis](1,1)--(2,1);\draw[axis](4,1)--(5,1);\draw[axis](1,-1)--(2,-1);\draw[axis](4,-1)--(3,-1);\node[above]at(2.5,1.5){频率 $\omega_s$};\node[below]at(2.5,-1.5){频率 $\omega_a$};
\end{tikzpicture}
\caption{测得的两频率分别对应中间弹簧不伸长与有伸长的模态。}\label{deep4:fig:d7-h3-q}
\end{figure}
\begin{solution}
图\ref{deep4:fig:d7-h3-q}的两模态满足 $m\omega_s^2=k$、$m\omega_a^2=k+2\kappa$，所以
\[
\boxed{k=m\omega_s^2,\qquad \kappa=\frac m2(\omega_a^2-\omega_s^2)>0.}
\]
振型 $(1,1)$ 与 $(1,-1)$ 都未按单位长度归一化。代入 $x_1=s+d$、$x_2=s-d$ 后，
\[
T=m(\dot s^2+\dot d^2)
=\frac12(2m)\dot s^2+\frac12(2m)\dot d^2.
\]
两个模态的惯性系数均为 $2m$。若误取单块质量m，每个模态的能量与总能量都会少算一半，而能量比仍保持不变。
初始同相、反相坐标为 $s_0=(A+B)/2$、$d_0=(A-B)/2$。相应等效质量均为 $2m$，故
\[
E_s=\frac m4\omega_s^2(A+B)^2,\qquad
E_a=\frac m4\omega_a^2(A-B)^2.
\]
\begin{figure}[H]\centering
\begin{tikzpicture}
\draw[axis](0,0)--(5.2,0)node[right]{$B/A$};\draw[axis](0,0)--(0,3.3)node[above]{模态能量};\draw[curve,domain=0:1,samples=70]plot({4.3*\x},{.65*(1+\x)^2});\draw[guide,domain=0:1,samples=70]plot({4.3*\x},{2.6*(1-\x)^2});\fill(1.433,1.156)circle(1.7pt);\draw[guide](1.433,1.156)--(2,2.2);\node[above]at(2.2,2.4){能量相等};\node[below]at(4.3,0){$1$};\node at(3.6,3){$\omega_a=2\omega_s$};
\end{tikzpicture}
\caption{同相能量随B增大而增加，反相能量减小；图示一组频率比。}\label{deep4:fig:d7-h3-s}
\end{figure}
由于 $0<B<A$，两式开正平方根即可得 $\omega_s(A+B)=\omega_a(A-B)$，因而
\[
\boxed{B=A\frac{\omega_a-\omega_s}{\omega_a+\omega_s}.}
\]
图\ref{deep4:fig:d7-h3-s}的交点唯一位于 $0<B/A<1$。模态能量相等并不要求模态位移系数相等，较高频模态在相同振幅下储存的能量较多。
\end{solution}


\begin{exercise}[失谐对拍振包络的抑制（选做）]\label{deep4:problem-exercise-6}\label{deep4:ex:d7-h6}
沿用例\ref{sup4:ex:d7-e5}的失谐双摆，取初始 $x_1=A,x_2=0$、两速度为零。记 $c=\kappa/m>0$，两频率为 $\omega_-<\omega_+$。求第二摆的位移及慢包络上界，并求包络上界降至 $A/2$ 时的失谐量。说明包络上界与实际位移最大值的区别。
\end{exercise}
\begin{figure}[H]\centering
\begin{tikzpicture}
\draw[axis](-1.8,0)--(5,0)node[right]{$x_1$};\draw[axis](0,-.5)--(0,3.5)node[above]{$x_2$};\draw[axis](0,0)--(2.8,1.5)node[right]{$\vct u_-$};\draw[axis](0,0)--(-1.5,2.8)node[above]{$\vct u_+$};\draw[curve](.8,0)arc(0:28:.8);\node at(1.1,.3){$\psi$};\fill(3.5,0)circle(2pt)node[below]{$(A,0)$};
\end{tikzpicture}
\caption{等质量时，可用相差直角的两个单位振型分解初始位移。}\label{deep4:fig:d7-h6-q}
\end{figure}
\begin{solution}
在图\ref{deep4:fig:d7-h6-q}中取 $\vct u_-=(\cos\psi,\sin\psi)$、$\vct u_+=(-\sin\psi,\cos\psi)$，其中 $0<\psi<\pi/2$。低频振型关系给出混合角的正弦和余弦：
\[
\cos2\psi=-\frac d{\sqrt{d^2+c^2}},\qquad
\sin2\psi=\frac c{\sqrt{d^2+c^2}}.
\]
它们等价于低频振型的 $\tan\psi=(d+\sqrt{d^2+c^2})/c$，使图\ref{deep4:fig:d7-h6-q}中的低频向量始终落在第一象限。仅用正弦关系会留下把两个摆球混合程度互换的歧义。
初态在两模态上的投影为 $A\cos\psi$、$-A\sin\psi$，所以
\[
x_2=A\cos\psi\sin\psi(\cos\omega_-t-\cos\omega_+t).
\]
令 $\bar\omega=(\omega_++\omega_-)/2$、$\Delta\omega=\omega_+-\omega_-$，得
\[
\boxed{x_2=\frac{Ac}{\sqrt{d^2+c^2}}\sin\bar\omega t\,
\sin\frac{\Delta\omega t}{2}.}
\]
\begin{figure}[H]\centering
\begin{tikzpicture}
\draw[axis](-3.4,0)--(3.6,0)node[right]{$d/c$};\draw[axis](0,0)--(0,2.7)node[above]{$A_{2,\rm env}^{\max}/A$};\draw[curve,domain=-3:3,samples=120]plot(\x,{2/sqrt(1+(\x)^2)});\draw[guide](-3,1)--(3,1);\draw[guide](1.732,0)--(1.732,1);\draw[guide](-1.732,0)--(-1.732,1);\node[below]at(1.732,0){$\sqrt3$};\node[below]at(-1.732,0){$-\sqrt3$};\node[left]at(0,2){$1$};\node[left]at(0,1){$1/2$};
\end{tikzpicture}
\caption{失谐使混合减弱，两侧对称地降低第二摆的包络上界。}\label{deep4:fig:d7-h6-s}
\end{figure}
由图\ref{deep4:fig:d7-h6-s}，包络上界为 $A_{2,\rm env}^{\max}=Ac/\sqrt{d^2+c^2}$。令它为 $A/2$，得到
\[
\boxed{|d|=\sqrt3c,}
\]
该失谐量还须满足 $|d|<\omega_c^2$。当 $|d|\gg c$，第二摆的响应明显减弱。两频率足够接近时，第二个正弦因子才描述慢包络；实际位移还乘有快速因子 $\sin\bar\omega t$。包络到达峰值时，快速因子未必恰为 $\pm1$，因此一拍内实际达到的最大位移可以小于这个上界。
\end{solution}


\begin{exercise}[中间质量改变后的三振子链]\label{deep4:problem-exercise-10}\label{deep4:ex:d7-h10}
例\ref{deep4:ex:d7-e9}的三振子链中，两外块质量仍为 $m$，中间块改为 $2m$，四根弹簧仍均为 $k$。求全部频率和振型，确定只对中间块施加冲量时不被激起的模态。
\end{exercise}
\begin{figure}[H]\centering
\begin{tikzpicture}
\foreach\x/\mass in{0/m,3/2m,6/m}{\draw[curve](\x,0)rectangle++(1,.7);\node at(\x+.5,.35){$\mass$};}\draw[spring](1,.35)--(3,.35);\draw[spring](4,.35)--(6,.35);\draw[spring](-1.5,.35)--(0,.35);\draw[spring](7,.35)--(8.5,.35);\draw[curve](-1.5,-.3)--(-1.5,1);\draw[curve](8.5,-.3)--(8.5,1);\draw[axis](3.5,1.2)--(5,1.2)node[right]{$J$};\node at(3.5,-.7){四弹簧均为 $k$};
\end{tikzpicture}
\caption{质量分布仍左右对称，中间质量改变不破坏镜面对称性。}\label{deep4:fig:d7-h10-q}
\end{figure}
\begin{solution}
图\ref{deep4:fig:d7-h10-q}中，反对称振型仍为 $(1,0,-1)$。中间块固定不动，两外块各受两根弹簧的回复力，因此 $\omega_a^2=2k/m$，与中间质量无关。

对称模态令 $x_1=x_3=x$、$x_2=y$，方程为
\[
m\ddot x+2kx-ky=0,\qquad
2m\ddot y+2ky-2kx=0.
\]
置 $\eta=m\omega^2/k$，得 $(2-\eta)x-y=0$、$-x+(1-\eta)y=0$。故
\[
\eta^2-3\eta+1=0,
\quad \boxed{\omega_\pm^2=\frac{k}{m}\frac{3\pm\sqrt5}{2}.}
\]
取外块振幅为1，对应振型为
\[
\vct u_-=(1,(1+\sqrt5)/2,1),\qquad
\vct u_+=(1,(1-\sqrt5)/2,1).
\]
\begin{figure}[H]\centering
\begin{tikzpicture}
\foreach\y/\txt in{2/低频,0/反对称,-2/高频}{\node[left]at(0,\y){\txt};\foreach\x in{1,3.5,6}{\fill(\x,\y)circle(2pt);}}
\draw[axis](1,2)--(1.8,2);\draw[axis](3.5,2)--(4.794,2);\draw[axis](6,2)--(6.8,2);
\draw[axis](1,0)--(1.8,0);\node[above]at(3.5,0){$u_2=0$};\draw[axis](6,0)--(5.2,0);
\draw[axis](1,-2)--(1.8,-2);\draw[axis](3.5,-2)--(3.006,-2);\draw[axis](6,-2)--(6.8,-2);
\end{tikzpicture}
\caption{改变中间质量后，两个对称模态的比例改变，反对称节点保留。}\label{deep4:fig:d7-h10-s}
\end{figure}
图\ref{deep4:fig:d7-h10-s}按频率从低到高排列：$(3-\sqrt5)/2<2<(3+\sqrt5)/2$。两个对称振型满足质量正交，因为令 $b_-=(1+\sqrt5)/2$、$b_+=(1-\sqrt5)/2$ 后，$b_-b_+=-1$，故
\[
\vct u_-^{\mathsf T}\operatorname{diag}(m,2m,m)\vct u_+
=2m+2mb_-b_+=0.
\]
中间冲量在反对称振型上的投影为零，该模态仍不出现。左右对称性保留了中间节点，中间质量的改变不影响这一选择规则。
把中间冲量记作 $\vct J=(0,J,0)^{\mathsf T}$，质量矩阵为 $M=\operatorname{diag}(m,2m,m)$。任一未归一化模态的初速度系数应由
\[
\dot q_j(0^+)=\frac{\vct u_j^{\mathsf T}\vct J}{\vct u_j^{\mathsf T}M\vct u_j}
\]
求出。图\ref{deep4:fig:d7-h10-s}中反对称振型的中间分量为零，分子随之为零；两个对称振型的中间分量非零，均会被激起。该判断同时用到了题设的零初位移，若初态已有反对称形变，该模态不会被一次中间冲量消除。
\end{solution}

